% Options for packages loaded elsewhere
\PassOptionsToPackage{unicode}{hyperref}
\PassOptionsToPackage{hyphens}{url}
\PassOptionsToPackage{dvipsnames,svgnames,x11names}{xcolor}
%
\documentclass[
  letterpaper,
  DIV=11,
  numbers=noendperiod]{scrartcl}

\usepackage{amsmath,amssymb}
\usepackage{iftex}
\ifPDFTeX
  \usepackage[T1]{fontenc}
  \usepackage[utf8]{inputenc}
  \usepackage{textcomp} % provide euro and other symbols
\else % if luatex or xetex
  \usepackage{unicode-math}
  \defaultfontfeatures{Scale=MatchLowercase}
  \defaultfontfeatures[\rmfamily]{Ligatures=TeX,Scale=1}
\fi
\usepackage{lmodern}
\ifPDFTeX\else  
    % xetex/luatex font selection
\fi
% Use upquote if available, for straight quotes in verbatim environments
\IfFileExists{upquote.sty}{\usepackage{upquote}}{}
\IfFileExists{microtype.sty}{% use microtype if available
  \usepackage[]{microtype}
  \UseMicrotypeSet[protrusion]{basicmath} % disable protrusion for tt fonts
}{}
\makeatletter
\@ifundefined{KOMAClassName}{% if non-KOMA class
  \IfFileExists{parskip.sty}{%
    \usepackage{parskip}
  }{% else
    \setlength{\parindent}{0pt}
    \setlength{\parskip}{6pt plus 2pt minus 1pt}}
}{% if KOMA class
  \KOMAoptions{parskip=half}}
\makeatother
\usepackage{xcolor}
\setlength{\emergencystretch}{3em} % prevent overfull lines
\setcounter{secnumdepth}{-\maxdimen} % remove section numbering
% Make \paragraph and \subparagraph free-standing
\makeatletter
\ifx\paragraph\undefined\else
  \let\oldparagraph\paragraph
  \renewcommand{\paragraph}{
    \@ifstar
      \xxxParagraphStar
      \xxxParagraphNoStar
  }
  \newcommand{\xxxParagraphStar}[1]{\oldparagraph*{#1}\mbox{}}
  \newcommand{\xxxParagraphNoStar}[1]{\oldparagraph{#1}\mbox{}}
\fi
\ifx\subparagraph\undefined\else
  \let\oldsubparagraph\subparagraph
  \renewcommand{\subparagraph}{
    \@ifstar
      \xxxSubParagraphStar
      \xxxSubParagraphNoStar
  }
  \newcommand{\xxxSubParagraphStar}[1]{\oldsubparagraph*{#1}\mbox{}}
  \newcommand{\xxxSubParagraphNoStar}[1]{\oldsubparagraph{#1}\mbox{}}
\fi
\makeatother


\providecommand{\tightlist}{%
  \setlength{\itemsep}{0pt}\setlength{\parskip}{0pt}}\usepackage{longtable,booktabs,array}
\usepackage{calc} % for calculating minipage widths
% Correct order of tables after \paragraph or \subparagraph
\usepackage{etoolbox}
\makeatletter
\patchcmd\longtable{\par}{\if@noskipsec\mbox{}\fi\par}{}{}
\makeatother
% Allow footnotes in longtable head/foot
\IfFileExists{footnotehyper.sty}{\usepackage{footnotehyper}}{\usepackage{footnote}}
\makesavenoteenv{longtable}
\usepackage{graphicx}
\makeatletter
\newsavebox\pandoc@box
\newcommand*\pandocbounded[1]{% scales image to fit in text height/width
  \sbox\pandoc@box{#1}%
  \Gscale@div\@tempa{\textheight}{\dimexpr\ht\pandoc@box+\dp\pandoc@box\relax}%
  \Gscale@div\@tempb{\linewidth}{\wd\pandoc@box}%
  \ifdim\@tempb\p@<\@tempa\p@\let\@tempa\@tempb\fi% select the smaller of both
  \ifdim\@tempa\p@<\p@\scalebox{\@tempa}{\usebox\pandoc@box}%
  \else\usebox{\pandoc@box}%
  \fi%
}
% Set default figure placement to htbp
\def\fps@figure{htbp}
\makeatother
% definitions for citeproc citations
\NewDocumentCommand\citeproctext{}{}
\NewDocumentCommand\citeproc{mm}{%
  \begingroup\def\citeproctext{#2}\cite{#1}\endgroup}
\makeatletter
 % allow citations to break across lines
 \let\@cite@ofmt\@firstofone
 % avoid brackets around text for \cite:
 \def\@biblabel#1{}
 \def\@cite#1#2{{#1\if@tempswa , #2\fi}}
\makeatother
\newlength{\cslhangindent}
\setlength{\cslhangindent}{1.5em}
\newlength{\csllabelwidth}
\setlength{\csllabelwidth}{3em}
\newenvironment{CSLReferences}[2] % #1 hanging-indent, #2 entry-spacing
 {\begin{list}{}{%
  \setlength{\itemindent}{0pt}
  \setlength{\leftmargin}{0pt}
  \setlength{\parsep}{0pt}
  % turn on hanging indent if param 1 is 1
  \ifodd #1
   \setlength{\leftmargin}{\cslhangindent}
   \setlength{\itemindent}{-1\cslhangindent}
  \fi
  % set entry spacing
  \setlength{\itemsep}{#2\baselineskip}}}
 {\end{list}}
\usepackage{calc}
\newcommand{\CSLBlock}[1]{\hfill\break\parbox[t]{\linewidth}{\strut\ignorespaces#1\strut}}
\newcommand{\CSLLeftMargin}[1]{\parbox[t]{\csllabelwidth}{\strut#1\strut}}
\newcommand{\CSLRightInline}[1]{\parbox[t]{\linewidth - \csllabelwidth}{\strut#1\strut}}
\newcommand{\CSLIndent}[1]{\hspace{\cslhangindent}#1}

\KOMAoption{captions}{tableheading}
\makeatletter
\@ifpackageloaded{caption}{}{\usepackage{caption}}
\AtBeginDocument{%
\ifdefined\contentsname
  \renewcommand*\contentsname{Table of contents}
\else
  \newcommand\contentsname{Table of contents}
\fi
\ifdefined\listfigurename
  \renewcommand*\listfigurename{List of Figures}
\else
  \newcommand\listfigurename{List of Figures}
\fi
\ifdefined\listtablename
  \renewcommand*\listtablename{List of Tables}
\else
  \newcommand\listtablename{List of Tables}
\fi
\ifdefined\figurename
  \renewcommand*\figurename{Figure}
\else
  \newcommand\figurename{Figure}
\fi
\ifdefined\tablename
  \renewcommand*\tablename{Table}
\else
  \newcommand\tablename{Table}
\fi
}
\@ifpackageloaded{float}{}{\usepackage{float}}
\floatstyle{ruled}
\@ifundefined{c@chapter}{\newfloat{codelisting}{h}{lop}}{\newfloat{codelisting}{h}{lop}[chapter]}
\floatname{codelisting}{Listing}
\newcommand*\listoflistings{\listof{codelisting}{List of Listings}}
\makeatother
\makeatletter
\makeatother
\makeatletter
\@ifpackageloaded{caption}{}{\usepackage{caption}}
\@ifpackageloaded{subcaption}{}{\usepackage{subcaption}}
\makeatother

\usepackage{bookmark}

\IfFileExists{xurl.sty}{\usepackage{xurl}}{} % add URL line breaks if available
\urlstyle{same} % disable monospaced font for URLs
\hypersetup{
  pdftitle={Reciprocal relationships, reverse causality, and temporal ordering: Testing theories with cross-lagged panel models},
  colorlinks=true,
  linkcolor={blue},
  filecolor={Maroon},
  citecolor={Blue},
  urlcolor={Blue},
  pdfcreator={LaTeX via pandoc}}


\title{Reciprocal relationships, reverse causality, and temporal
ordering: Testing theories with cross-lagged panel models}
\author{}
\date{}

\begin{document}
\maketitle


Reciprocal causal relationships are a common feature of criminological
theories. For example, stable employment may reduce offending while
offending may lead to unemployment, and perceived disorder may increase
fear of crime while fear of crime may increase sensitivity to signs of
disorder. When multiple observations over time are available,
cross-lagged panel models are commonly used to estimate these reciprocal
effects. Yet this is often done without careful attention to how they
map on to the theoretical process they are meant to capture or whether
key assumptions of the models are satisfied. This may result in
estimates that are not substantively meaningful or are biased or even
reversed in sign. Reciprocal relationships also pose challenges for
causal assumptions based on graphical tools; theories that posit
reciprocal causation often rely on underlying macro--micro mechanisms
not explicitly represented in empirical models. We provide guidance on
how to align theory, model specification, and choice of estimator,
illustrated with two motivating examples: the relationship between
employment and offending and the relationship between perceived disorder
and fear of crime. We distinguish theoretical reciprocality---in which
feedback processes are of substantive interest---from reverse
causality---in which reciprocal processes are a nuisance threatening
identification of an effect of interest---and discuss three common
issues for cross-lagged panel models: unobserved time-stable
heterogeneity, improperly specified temporal order, and insufficient
variation over time. We conclude by reminding researchers of the
importance of aligning methods with theory in longitudinal research.

\subsection{Declarations}\label{declarations}

The authors did not receive support from any organization for the
submitted work.

\section{Introduction}\label{introduction}

Reciprocal relationships and reverse causality are a common feature of
criminological theories. In some cases, theories suggest that two
variables are reciprocally related, reproducing each other over time in
a cycle. For example, labeling theory suggests individuals labeled as
criminals are more likely to engage in criminal activity, which then
contributes to reinforcing their labels as criminals, forming a cyclic
pattern of criminal reproduction
(\citeproc{ref-becker1963outsiders}{Becker 1963}). In other cases,
competing theories suggest distinct explanations regarding the causal
direction of an established association. For example, the broken windows
thesis posits social and physical disorder in the neighborhood increases
people's fear of crime, whereas socio-psychological theories suggest
that fear of crime contributes to heightened perceptions of neighborhood
disorder (\citeproc{ref-sampson2004seeing}{Sampson and Raudenbush 2004};
\citeproc{ref-jackson2004experience}{Jackson 2004}). In still other
cases, theories suggest a causal effect of a treatment on an outcome,
while also cautioning that the outcome can affect the treatment over
time. For example, developmental and life-course criminology suggests
that stable employment reduces criminal behavior, whereas criminal
behavior can also lead to job loss or unemployability
(\citeproc{ref-ramakers2017not}{Ramakers et al. 2017};
\citeproc{ref-crutchfield2014get}{Crutchfield 2014}). Across all these
examples, reciprocality---the possibility that two variables influence
each other over time---is at the heart of the theoretical process.

When repeated observations over time are available, the cross-lagged
panel model (CLPM) is one of the most commonly used tools for studying
these reciprocal processes. In its simplest form, a CLPM regresses each
variable on its own lagged value and on the lagged value of the other
variable, producing autoregressive and cross-lagged parameter estimates.
A credible non-zero cross-lagged coefficient is typically interpreted as
evidence that one variable influences the other over time, net of their
respective stabilities. Because both cross-lagged paths can be estimated
simultaneously, the CLPM appears to offer a convenient way to assess
reciprocality or to adjudicate between competing causal directions---and
it has been applied to a wide range of topics in criminology and beyond
(e.g., \citeproc{ref-matsueda1998dynamics}{Matsueda and Anderson 1998};
\citeproc{ref-brunton2011untangling}{Brunton-Smith 2011};
\citeproc{ref-pina2020reassessing}{Pina-Sánchez and Brunton-Smith
2020}).

Yet CLPMs are frequently used without careful attention to the
theoretical process they are meant to capture. In particular, studies
sometimes fail to consider the role of time when specifying temporal
processes. Causal processes unfold at particular speeds---for example,
while employment may constrain offending on a daily basis simply by
occupying one's time, the reputational consequences of a criminal record
may take months to translate into job loss. These temporal dynamics
should inform not only modeling strategies but, more importantly, the
definition of the causal quantity of interest---the \emph{estimand},
which is often left implicit, determined not by theory but by the
structure of the available data. When the starting point is not a
clearly specified causal process, one risks arriving at an estimator
that produces a correct answer to the wrong question. In addition, the
apparent simplicity of the CLPM belies a number of assumptions that must
be satisfied for its estimates to be valid. These include proper
accounting for unobserved heterogeneity that might confound the
estimated relationships, correctly specified temporal lags between
exposure and outcome, and sufficient within-unit variation over time to
detect cross-lagged effects. Failure to satisfy these assumptions can
produce severe problems, including spurious associations, imprecise
estimates, and parameter estimates that are biased or even reversed in
direction.

Our goal in this paper is to provide practical guidance for researchers
studying reciprocal relationships or dealing with reverse causality
using panel data. We do not develop new estimators or derive new formal
results, but we synthesize several recent methodological contributions
that have been developed by different disciplinary traditions (e.g.,
\citeproc{ref-allison2017maximum}{Allison, Williams, and Moral-Benito
2017}; \citeproc{ref-hamaker2015critique}{Hamaker, Kuiper, and Grasman
2015}; \citeproc{ref-zyphur2020data}{Zyphur et al. 2020}) in a way that
is accessible to applied researchers in criminology. First, we highlight
the distinction between \emph{theoretical reciprocality}, when feedback
processes are the core of the theories being tested, and \emph{reverse
causality}, when reciprocal effects are a nuisance of the theory being
tested. We focus on three key problems that we believe are the most
consequential for applied research in developmental and life-course
criminology: (a) unobserved heterogeneity, (b) improperly specified
temporal lags, and (c) insufficient inter-temporal variation. These are
not the only issues that can arise with CLPMs, but they are the ones
that researchers will most frequently encounter and that can have the
most severe consequences for inference. For each, we explain why it
matters, how it arises in practice, and what solutions are available.

We ground our discussion in two motivating examples that represent the
two broad reasons one might turn to a CLPM. The first is the
relationship between employment and criminal offending---a staple of
life-course criminology (\citeproc{ref-sampson1995crime}{Sampson and
Laub 1995}; \citeproc{ref-laub2006shared}{Laub and Sampson 2006}). Here,
the primary research question concerns the causal effect of employment
on crime, but reverse causality (i.e., crime involvement leading to
unemployment) threatens to bias estimates of this effect. In this case,
reciprocality is a nuisance to be controlled for. The second is the
relationship between perceived neighborhood disorder and fear of
crime---a central topic for neighborhood researchers
(\citeproc{ref-sampson2004seeing}{Sampson and Raudenbush 2004};
\citeproc{ref-obrien2019looking}{O'Brien, Farrell, and Welsh 2019}).
Here, competing theories propose opposite causal directions, and the
reciprocal relationship itself is the substantive question. We use these
examples to illustrate how the same general class of models can serve
different research purposes and how the distinction between reverse
causality as a nuisance and theoretical reciprocality as a substantive
question should also shape model specification.

The paper proceeds as follows. We first discuss estimands and the role
of time, highlighting what we mean by reciprocal effects in the context
of theory testing and the role of unexamined mechanisms operating at a
lower, more temporally disaggregated level than observed measures. We
then present our two motivating examples in detail, with a focus on the
distinction between reverse causality and theoretical reciprocality. We
then develop the statistical framework and discuss why we recommend
using directed acyclic graphs (DAGs) to represent the theoretical model
and structural equation models (SEMs) to represent the statistical
model, while being attentive to where these two representations
necessarily diverge. The main body of the paper addresses the three main
issues---unobserved heterogeneity, temporal order, and insufficient
inter-temporal variation---as well as solutions for each of them. We
conclude with general recommendations for researchers studying
reciprocal relationships with panel data.

\section{Estimands and the role of
time}\label{estimands-and-the-role-of-time}

An estimand is the quantitative definition of the estimate of interest,
such as a causal effect (see \citeproc{ref-lundberg2021your}{Lundberg,
Johnson, and Stewart 2021}). In a counterfactual framework, an estimand
may be defined as a difference in potential outcomes
(\citeproc{ref-morgan2014counterfactuals}{Morgan and Winship 2014}),
e.g., ``how much higher is the probability that an individual will
offend in a given month had they become unemployed in the prior month?''
Estimands do not exist out in the world waiting to be found; they are
chosen by researchers. We defer to other sources for a full explanation
of choosing and using estimands (e.g.,
\citeproc{ref-lundberg2021your}{Lundberg, Johnson, and Stewart 2021}).
For the present work, what is important is recognizing that there is no
``true causal timing'' waiting to be discovered because causal timing is
a choice a researcher makes when defining their estimand of interest.

To illustrate, consider this: what is the true causal timing of drinking
an additional pint of beer on one's perceived wellbeing? What is this
effect at 20 seconds, 20 minutes, 20 hours, or 20 years? What about the
effect of an additional pint per day after 10 years? In this context,
the reader may know intimately which exposure and which lags between
measurement of exposure and outcome correspond to detectable and
substantively interesting causal effects. Strong theories should be
specific about the role of time and thus able to inform what sort of
causal timings---whether delay between exposure and outcome or
aggregations of exposures---are of substantive interest
(\citeproc{ref-hamaker2026}{Hamaker 2026};
\citeproc{ref-nivette2026}{Nivette and Eisner 2026}). Causal processes
are often (but not always) continuous but researchers discretize time
for practical utility, e.g., we know drinking a pint of beer causes
one's blood alcohol level to ramp up then back down but someone may be
most interested in knowing how they will feel 30 minutes from now or
when they wake up in 12 hours---or how it will impact their health in
ten years if they make it a regular habit.

Often the estimates produced in longitudinal research do not correspond
to any theoretically-motivated or substantively useful estimand. While a
relatively common specification in longitudinal research, consider when
we would be interested in the (lagged) effect of proportion of time
employed last year on crimes committed this year. Does this estimand
make sense if one believes employment reduces crime by keeping people
busy? Implicit estimands such as these are often chosen by starting with
a methods concern (e.g., reverse causality) and working backwards, so
one arrives at a model that gives a correct answer to the wrong
question.

Relatedly, researchers should recognize that, because causal effects
occur over time, there are no simultaneous non-recursive causal
relationships (\citeproc{ref-strotz1960recursive}{Strotz and Wold
1960}). When researchers posit non-recursive relationships, they are
nearly always conducting analyses at a higher-level (i.e., macro) than
where the process of interest is taking place (i.e., micro). For
example, while economics depicts supply and demand as simultaneously
determined, this relationship becomes recursive if we make explicit the
underlying micro-macro model in which individual actors on both sides
adjust to each others' behavior over time to produce the equilibrium
price. Abstraction and aggregation to higher levels of analysis is both
useful and ubiquitous throughout the sciences, and formal guidelines
exist for when macro-level causal models will accurately represent (or
translate) the outcomes of micro-level causal models (e.g.,
\citeproc{ref-rubenstein2017causal}{Rubenstein et al. 2017}). These
abstractions are often forced on researchers by limitations in
measurement. Long time lags in longitudinal research are a special case
of this micro-macro issue in which faster-operating micro processes
(e.g., daily processes) are modeled at a longer macro time scale (e.g.,
using waves of an annual survey).

Consequently, researchers should make an argument that the timing of the
effect in the estimand they are studying is useful and can be studied
with the data they have. While we focus here on longitudinal research,
this applies just as much to cross-sectional research, except there
researchers have even less control over what timings they may examine.
Even provided otherwise identical estimands, discrepant results between
studies may be due to different choices of timing
(\citeproc{ref-gollob1987taking}{Gollob and Reichardt 1987}; see also
\citeproc{ref-vaisey2017you}{Vaisey and Miles 2017}); the effect of an
additional pint on perceived wellbeing will be quite different measured
in one hour versus in the following morning.

\section{Two motivating examples}\label{two-motivating-examples}

To illustrate our points, we consider two canonical longitudinal
relationships in criminology. The first is the link between employment
and offending---a focus of classic life‐course research (e.g.,
\citeproc{ref-sampson1995crime}{Sampson and Laub 1995};
\citeproc{ref-laub2006shared}{Laub and Sampson 2006};
\citeproc{ref-crutchfield2014get}{Crutchfield 2014}). The second is the
link between residents' perceptions of neighborhood disorder and their
fear of crime---a central topic for neighborhood scholars
(\citeproc{ref-sampson2004seeing}{Sampson and Raudenbush 2004};
\citeproc{ref-obrien2019looking}{O'Brien, Farrell, and Welsh 2019}). In
the employment--crime example, classical criminological theory and
policy emphasize the effect of employment on future offending, whereas
in the fear--disorder example the direction of causation is itself
debated.

\subsection{Employment and crime}\label{employment-and-crime}

Life‐course studies routinely find that obtaining stable employment
reduces criminal behavior (\citeproc{ref-ramakers2017not}{Ramakers et
al. 2017}; \citeproc{ref-lageson2012work}{Lageson and Uggen 2012}). The
logic is straightforward. In the short run, being employed occupies
one's time and builds conventional social bonds, which constrains
opportunities and incentives for offending
(\citeproc{ref-hirschi1969causes}{Hirschi 1969}). In the long run, it is
not employment per se but stable, committed work that enhances social
control and thereby reduces crime
(\citeproc{ref-sampson1995crime}{Sampson and Laub 1995};
\citeproc{ref-laub2006shared}{Laub and Sampson 2006}). In other words,
the accumulation of conventional ties and responsibilities that are
often associated with stable employment provides an identity that deters
future offending (\citeproc{ref-ramakers2017not}{Ramakers et al. 2017}).
For example, longitudinal reviews show that employed individuals have
significantly lower recidivism rates than the unemployed
(\citeproc{ref-uggen2008have}{Uggen and Wakefield 2008}; although see
\citeproc{ref-crutchfield2014get}{Crutchfield 2014}), and experimental
job‐program evaluations have found that random assignment to job
training can yield lower subsequent crime
(\citeproc{ref-redcross2011more}{Redcross et al. 2011}). The effect of
work on offending is of particular interest for intervention research,
because it is easier for policy to manipulate employment than other
candidate causes of crime such as deviant peers or neighborhood
conditions and to target it at high-risk populations
(\citeproc{ref-uggen2008have}{Uggen and Wakefield 2008}).

Crucially, however, the reverse causal path is also possible. After all,
crime may serve as a substitute for conventional employment
(\citeproc{ref-fagan1999}{Fagan and Freeman 1999}) and arrests or
incarceration can lead to job loss, inhibit human capital accumulation,
and make getting a job difficult
(\citeproc{ref-crutchfield2014get}{Crutchfield 2014};
\citeproc{ref-chalfin2011}{Chalfin and Raphael 2011};
\citeproc{ref-pager2003}{Pager 2003}). Therefore, if interested in
estimating the effects of employment on crime, one needs to account for
the effects of crime on employment---i.e., one needs to adjust for
\emph{reverse causality}. Here we adopt the interventionist perspective
and focus on the the effect of employment on crime, treating effects of
crime on employment as reverse causality that needs to be controlled
for. This is, of course, a choice. For example, a researcher interested
in employment and offending in a long term desistance process (e.g.,
\citeproc{ref-giordano2002}{Giordano, Cernkovich, and Rudolph 2002}) or
in the fluidity of engagement in conventional work and crime
(\citeproc{ref-fagan1999}{Fagan and Freeman 1999}) might instead
approach this as a \emph{theoretical reciprocality} problem. As with
estimands before, how to approach reciprocality is a choice by the
researcher that should be motivated by theory and the research question.

\subsection{Fear of crime and perceived
disorder}\label{fear-of-crime-and-perceived-disorder}

By contrast, the relationship between perceived disorder and fear of
crime is often viewed as reciprocal
(\citeproc{ref-brunton2011untangling}{Brunton-Smith 2011}). Major
perspectives in criminology posit that visible disorder in the
environment raises fear of crime (\citeproc{ref-taylor2001}{Taylor
2001})---for example, that disorderly cues signal that an area is in
decline and unable to control deviance
(\citeproc{ref-kelling1982broken}{Wilson and Kelling 1982}). While they
do not account for potentially biasing reverse causality, a review by
O'Brien et al. (\citeproc{ref-obrien2019looking}{2019}) argues there is
stronger evidence for an effect of perceived disorder on fear of crime
of than of objective disorder. Others emphasize the opposite
interpretation, finding that individuals with high levels of fear of
crime are sensitized to local disorder, causing them to perceive and
interpret ambiguous cues as physical and social disorder
(\citeproc{ref-sampson2004seeing}{Sampson and Raudenbush 2004};
\citeproc{ref-jackson2004experience}{Jackson 2004})---a relationship
which may be related to the stronger observed effects of perceived
versus objective disorder on fear. In practice, both processes may
occur. Given the theoretical ambiguity and the possibility for
reciprocal relationships over time, this example does not have clearly
defined dependent and independent variables: both fear of crime and
perceived disorder may be seen as outcomes as well as predictors of
interest. The substantive interest is therefore the reciprocality
itself. In this case, one can use panel models either to adjudicate
between competing theories or to discover actual reciprocal effects over
time. For example, Brunton-Smith
(\citeproc{ref-brunton2011untangling}{2011}) found evidence only for an
effect of disorder on fear in the next year (and not vice versa) using
cross-lagged panel model with yearly survey data---but could not test
for faster operating (i.e., within-year) effects which could occur if
fear and disorder perceptions result from rapid updating processes.

\section{Theoretical and statistical
models}\label{theoretical-and-statistical-models}

These examples are contrasted by how they are framed theoretically. In
the employment--crime case, the cross‐lagged panel model is used to
account for reverse causality (crime \(\rightarrow\) employment) and
isolate the main effect of interest (employment \(\rightarrow\) crime).
In the fear--disorder case, both cross-lagged paths (disorder
\(\rightarrow\) fear; fear \(\rightarrow\) disorder) are of substantive
interest, reflecting genuine theoretical ambiguity. Which direction we
emphasize is dictated by the research question. In the first example we
treat the reverse path as a nuisance (i.e., reverse causality), whereas
in the second we treat it as part of a mutually reciprocal process to be
tested (i.e., theoretical reciprocality).

Good empirical research starts with a strong theoretical foundation. It
is from a clearly defined theoretical framework that observable
implications can be deduced, and it is those observable implications
that can then be contrasted against empirical data
(\citeproc{ref-king1994designing}{King, Keohane, and Verba 1994};
\citeproc{ref-mayo2018statistical}{Mayo 2018}). Theory-testing empirical
research is therefore conditional on a thorough and transparent
theoretical model (\citeproc{ref-nagin2019real}{Nagin and Sampson
2019}). Given that most criminological theories presume causal
processes, the first step should normally involve deriving a recursive
model that accurately represents the theorized relationships---e.g., the
theoretical model could be represented by a directed acyclic graph
(DAG). DAGs can formally represent a nonparametric structural equation
model, avoiding functional-form and distributional assumptions
(\citeproc{ref-pearl2009causality}{Pearl 2009}), and are increasingly
used to portray premised causal relationships in criminological theory
(\citeproc{ref-nivette2025assessing}{Nivette and Vegt 2025};
\citeproc{ref-oliveira2024aggressive}{Oliveira 2024};
\citeproc{ref-lanfear2024promise}{Lanfear and Kirk 2024}). The DAG
should be comprehensive and include all relevant theorized pathways,
even those involving unobserved variables.

Once the theoretical model is specified, it should guide the selection
of an appropriate estimator. It is important to distinguish the
theoretical model from the statistical model: while the latter should
approximate the former as much as possible, the statistical model will
unavoidably have to handle imprecision and uncertainty that are bound to
exist when dealing with empirical data, such as distributional and
parametric assumptions and unanalyzed covariances due to theoretical
ambiguities. When dealing with situations involving potential reciprocal
relationships and/or reverse causality, we propose leveraging DAGs for
theoretical models and drawing on the structural equation modeling (SEM)
framework for estimation (\citeproc{ref-bollen2013eight}{Bollen and
Pearl 2013}). A clear theoretical model is a necessary prerequisite for
specifying an estimator. While both diagrams will not be exactly the
same, we suggest that the premised DAG inform some elements of the SEM.
Specifically, unidirectional arrows in the structural model should be
recursive, as per the theoretical model.\footnote{SEMs include a
  measurement component, which specifies measurement models linking
  manifest indicators and latent variables, and a structural component,
  which specifies the relationships between latent variables. For now,
  we are exclusively referring to the structural component of an SEM.
  Those include bi-directional arrows, which depict the estimation of a
  covariance between two variables (or their error terms if they are
  endogenous), and unidirectional arrows, which indicate the estimation
  of directional associations based on regression modeling.} For
example, as depicted in Figure~\ref{fig-classic-clpm}, a key
consideration in this process is the timing of effects: theorized
effects can be contemporaneous (i.e., occur quickly) or lagged (i.e.,
unfold more gradually over time).

\begin{figure}

\centering{

\pandocbounded{\includegraphics[keepaspectratio]{./jdlcc_draft_deidentified_files/figure-pdf/fig-classic-clpm-1.pdf}}

}

\caption{\label{fig-classic-clpm}Directed acyclic graph representation
of cross-lagged panel models with (a) lagged effects only or (b)
contemporaneous effects as well}

\end{figure}%

Yet, while the theoretical model displays a non-parametric
representation of theorized causal relationships, the SEM can include
bi-directional arrows representing empirically observed covariances
between variables not explained by the premised DAG. Panel (a) in
Figure~\ref{fig-sem-clpm} depicts a statistical model corresponding to
the theoretical model in panel (a) of Figure~\ref{fig-classic-clpm},
where the reciprocal effects of interest between \(X\) and \(Y\) occur
only across periods but we additionally include unanalyzed (or nuisance)
covariances between them to account for contemporaneous causal effects
of ambiguous direction. In contrast, panel (b) in
Figure~\ref{fig-sem-clpm}, which corresponds to the theoretical model in
panel (b) of Figure~\ref{fig-classic-clpm}, features contemporaneous
structural paths \(X_t \rightarrow Y_t\). This means there are no longer
any unanalyzed covariances, so the statistical model is nearly identical
to the DAG.\footnote{Provided an instrument for \(X_t\), one could add a
  covariance between \(\zeta_t\) and \(\epsilon_t\) to adjust for
  contemporaneous reverse causality.} Note here \(Y_1\) is now
endogenous (and has an error term) while it was exogenous (no error
term) in panel (a). The models in both panels assume no serial
correlation as indicated by the absence of paths between the error terms
in different time periods. While, unlike a DAG, an SEM permits
contemporaneous reciprocal causal effects, we suggest that the necessary
non-recursive simultaneous equations be used only as a last resort,
given their often tenuous theoretical grounding and in particular strong
assumptions required for identification
(\citeproc{ref-wooldridge2010econometric}{Wooldridge 2010}).\footnote{Plausible
  instrumental variables are notoriously difficult to find and
  non-recursive simultaneous equations require instruments for
  \emph{each} endogenous variable of interest.}

\begin{figure}

\centering{

\pandocbounded{\includegraphics[keepaspectratio]{./jdlcc_draft_deidentified_files/figure-pdf/fig-sem-clpm-1.pdf}}

}

\caption{\label{fig-sem-clpm}SEM path diagram representations of
three-wave CLPM. Panel (a) features only lagged effects but permits
ambiguous contemporaneous covariances between \(X\) and \(Y\). Panel (b)
features contemporaneous effects of \(X\) on \(Y\).}

\end{figure}%

\subsection{When is a CLPM needed?}\label{when-is-a-clpm-needed}

A cross-lagged panel model is only needed when reciprocal effects over
time are presumed in the theoretical model. When the researcher can
theoretically assume no effects of the dependent variable on the
independent variable over time, other statistical models will be more
appropriate than a cross-lagged panel model (see
\citeproc{ref-imai2019should}{Imai and Kim 2019}). For example, research
focused on the effects of a single-event exposure---including the role
of turning points commonly studied in life-course criminology such as
the long-term effects of school drop-out or marriage on offending
(\citeproc{ref-sampson2016turning}{Sampson and Laub 2016};
\citeproc{ref-nguyen2018measurement}{Nguyen and Loughran 2018})---would
not benefit from cross-lagged panel models. Instead, reciprocal effects
over time are most often a threat when the exposure is a continuous
process occurring repeatedly. There are two situations in which
reciprocal effects could emerge from theoretical models, thus
highlighting the need of cross-lagged panel models. We refer to them as
\emph{reverse causality} and \emph{theoretical reciprocality}.

Reverse causality, on the one hand, refers to situations in which
researchers are primarily interested in assessing the effects of an
exposure on an outcome but are concerned about the possibility that the
outcome also influences future exposures. From our empirical example,
assessing the effects of employment on crime involvement is of
particular theoretical and policy relevance
(\citeproc{ref-sampson1995crime}{Sampson and Laub 1995};
\citeproc{ref-crutchfield2014get}{Crutchfield 2014}); in order to assess
that relationship, however, one needs to take into account that
offending behavior may also have an effect on future employment. This is
a classic endogeneity problem, in which not accounting for reverse
causality introduces bias. This has long been a concern of
econometricians---e.g., the Arellano-Bond estimator, which is
essentially a cross-lagged panel model, was developed specifically to
address endogeneity in dynamic economic relationships
(\citeproc{ref-arellano1991some}{Arellano and Bond 1991};
\citeproc{ref-moral2019dynamic}{Moral-Benito, Allison, and Williams
2019}). When dealing with reverse causality, researchers have a
predefined exposure and an outcome of interest and simply aim to adjust
away the effects of the outcome on the exposure over time.

What we call theoretical reciprocality, on the other hand, refers to
situations in which the directionality of the effects is ambiguous or of
substantive interest. This could be due to, for instance, competing
theories suggesting opposing pathways between two variables
(\citeproc{ref-sturgis2004using}{Sturgis et al. 2004}). From our
empirical example, some researchers argue that perceived disorder leads
to increases in fear of crime (\citeproc{ref-kelling1982broken}{Wilson
and Kelling 1982}) whereas others approach this relationship from the
opposite direction, arguing that fear of crime leads local residents to
view disorderly signs in the neighborhood as problematic
(\citeproc{ref-sampson2004seeing}{Sampson and Raudenbush 2004};
\citeproc{ref-jackson2004experience}{Jackson 2004}). This theoretical
ambiguity implies that reciprocality is an empirical problem---a
question that can be adjudicated by theory testers
(\citeproc{ref-brunton2011untangling}{Brunton-Smith 2011}). When dealing
with theoretical reciprocality, researchers often do not have a
predefined exposure and an outcome of interest and instead aim to assess
both directions.

In both cases, the theoretical model will highlight the threat of
reciprocal relationships. A DAG will not distinguish between reciprocal
effects due to theoretical reciprocality or due to reverse causality;
this distinction is purely conceptual and should be dictated by the
research question. Yet, we maintain that this is an important
distinction. While in some cases the empirical solution will be similar,
in other situations dealing with theoretical reciprocality or reverse
causality may lead to different statistical estimators---we expand on
this point below.

\subsection{Panel Data}\label{panel-data}

For the discussion that follows, let us assume that in both example
cases we have panel data with three yearly waves of observation and
continuous measures of our concepts. In our first example---the effect
of employment on offending and attendant reverse causality---employment
is measured as the proportion of weeks spent employed in the last year
and offending is a self-reported count of offenses committed in the past
year. Both measures are taken at survey time but describe conditions or
events occurring over the span of the year prior to each survey wave. In
our second example---the theoretical reciprocal effects of perceived
disorder and fear of crime---both perceived disorder and fear of crime
are measured at survey time. That is, we observe respondents'
``current'' perceptions of disorder in their neighborhood and their
``current'' fear of crime in that neighborhood. To simplify examples, we
will assume this is measured using a validated battery of indicators
that produce a proportional score from 0 (no perceived disorder or fear)
to 1 (maximum possible disorder or fear).

Theory is recursive and operates at some defined pace---often discrete
but sometimes continuous---but practical constraints of measurement
commonly result in data collected in waves. The widths of time intervals
between waves are usually determined by pragmatic considerations rather
than theory, and thus the timing of observations rarely corresponds to
the timing we expect in our theories. Often timing is inconsistent
across variables as well, e.g., measures describing activities in the
past year sit alongside measures of present perceptions. More generally,
we can imagine the discordance between our theoretical model and our
empirical data collection as a missing data or latent variable problem
as illustrated in Figure~\ref{fig-time-dag-unobs}. As noted, in the real
world, time is continuous, so one might imagine there exists an infinite
number of these ``missing waves'' between observations. Whether explicit
or not, researchers always make assumptions about the way variables and
relationships operate within these unobserved intervals. Theory should
and sometimes even does provide a guide to what assumptions are
reasonable to make (e.g., see \citeproc{ref-nivette2026}{Nivette and
Eisner 2026}).

\begin{figure}

\centering{

\pandocbounded{\includegraphics[keepaspectratio]{./jdlcc_draft_deidentified_files/figure-pdf/fig-time-dag-unobs-1.pdf}}

}

\caption{\label{fig-time-dag-unobs}A theoretical model of cross-lagged
effects where observations do not correspond to the timing of the
underlying causal process. Circled variables are latent.}

\end{figure}%

When observation periods correspond in some way to the speed at which
causal processes are believed to operate, theory may imply that one can
rule out reciprocal effects or reverse causality. For example, if
employment were measured as ``Are you presently employed?'', one could
assume employment can be affected by past-year offending but employment
cannot affect past-year offending. In neither of our present examples,
however, do our empirical observation periods align with the theoretical
causal process. Here we focus on cases where the processes in which we
are interested occur more quickly than data are
recorded\footnote{Mulder, Voelkle, and Hamaker
  (\citeproc{ref-mulder2025time}{2025}) also note the possibility of
  ``low process coverage'' that occurs when observation intervals that
  are too long combine with time aggregated variables to make it so an
  insufficient amount of the underlying continuous process is captured
  by the data.}---that is, we must model a micro-level process using a
macro-level aggregation. The reverse case---when recording occurs more
quickly than the causal process---is easily addressed by aggregation or
temporal lags.\footnote{This is similar to using lags of variables from
  neighbors in spatial models to address the possibility that processes
  do not perfectly map on to the chosen geographic units.}

While we treat time as discrete in both theory and analyses---as this is
the dominant approach in the field of criminology---it is also possible
to model cross-lagged relationships using continuous time dynamic models
(see \citeproc{ref-voelkle2018role}{Voelkle et al. 2018}). Discrete-time
CLPMs make an implicit assumption that the chosen observation interval
is the relevant temporal scale for the causal process: the estimated
lagged coefficients are specific to that interval and would generally
differ had observations been collected on a different schedule
(\citeproc{ref-gollob1987taking}{Gollob and Reichardt 1987}).
Continuous-time dynamic structural equation models
(\citeproc{ref-voelkle2018role}{Voelkle et al. 2018};
\citeproc{ref-driver2017continuous}{Driver, Oud, and Voelkle 2017})
instead parameterize the underlying continuous-time dynamics, from which
the implied discrete-time relationships at any interval can be derived.
This makes effect estimates comparable across studies with different
observation intervals and allows researchers to project effects onto
intervals not directly observed. However, the continuous-time framework
does not remove the need for adequate temporal resolution. Identifying
the underlying continuous-time process from a sparse discrete-time panel
is generally not uniquely solvable without either frequent observations
or strong functional-form assumptions about how processes unfold between
waves. Continuous-time models are therefore most useful when researchers
have either fine-grained data (e.g., experience sampling) or strong
theoretical commitments about within-interval dynamics.\footnote{See
  Nivette \& Eisner (\citeproc{ref-nivette2026}{2026}) for an extended
  discussion of the relationship between continuous time processes and
  discrete observations in developmental and life course research.}

\subsubsection{Temporality and composite
variables}\label{temporality-and-composite-variables}

Before proceeding to the statistical model, it is necessary to discuss
problems that can result from measurement strategies common in panel
studies. For this purpose, we will distinguish between two types of
measure, which we will call \emph{instantaneous measures} and
\emph{period measures}. An instantaneous measure is one where the
measured values correspond to some single instant in time (or at least
an exceedingly short period). Our perceived disorder and fear of crime
variables are examples of instantaneous measures, as they would
typically be measured using respondent perceptions at survey time. This
situation corresponds roughly to that depicted in
Figure~\ref{fig-time-dag-unobs} and can be addressed with conventional
approaches to causal inference with DAGs.\footnote{Note that the
  theoretical causal variable of interest may instead be this variable's
  value in some prior period; the available instantaneous measure thus
  represents a proxy for the past variable. The measure may be a poor
  proxy if the underlying variable varies significantly over time or is
  subject to time-varying confounding.}

In contrast, a period measure is one where the measured value
corresponds to some deterministic sum (or other aggregate) calculated
over a specified period of time. A count or variety score of offenses
committed over the past year and proportion of time spent employed in
the past year are both examples of period measures. When our observation
periods do not correspond to the timing of our underlying theoretical
causal model, these period measures become what we will call
\emph{composite period variables}. They are a composite because the
measured value is a deterministic sum of values that would be obtained
by dividing the observation period up into smaller intervals and
aggregating across them.\footnote{See also Mulder, Voelkle, and Hamaker
  (\citeproc{ref-mulder2025time}{2025}) for a discussion of ``time
  aggregation'' where the measurement represents some average over the
  time period \emph{as recalled by respondents}.} This situation,
appearing in Figure~\ref{fig-app-latent}, produces challenges that are
not well recognized in the literature and sometimes difficult to address
(see \citeproc{ref-berrie2025depicting}{Berrie et al. 2025}).

\begin{figure}

\centering{

\pandocbounded{\includegraphics[keepaspectratio]{./jdlcc_draft_deidentified_files/figure-pdf/fig-app-latent-1.pdf}}

}

\caption{\label{fig-app-latent}Reciprocality due to composite period
variable. Circled variables are latent. Double-arrows indicate observed
variables are composite variables determined by their parents}

\end{figure}%

Figure~\ref{fig-app-latent} illustrates a case where we are interested
in the relationship between two period variables \(X\) and \(Y\). The
recall period for data collection encompasses 1 unit of time (e.g., an
entire year) but the causal process relating the two period variables
operates over smaller intervals (e.g., six months) denoted \(A\) and
\(B\). This is represented by cross-lagged effects between latent \(X\)
and \(Y\) variables occurring across each smaller interval. Because data
were only collected at each observation time, the measured values
represent a composite period variable combining the values from periods
\(A\) and \(B\). For example, if \(Y_1\) is the number of self-reported
criminal offenses in a year, it would be the simple sum of the offenses
in each six month period \(Y_{1A} + Y_{1B}\). While we may know this
sum, without additional information, we are uncertain of the value of
either \(Y_{1A}\) or \(Y_{1B}\) (unless \(Y_1=0\)). Borrowing the
notation of Berrie et al. (\citeproc{ref-berrie2025depicting}{2025}),
the double-arrows indicate that \(Y_t\) and \(X_t\) are
\emph{deterministic} variables, i.e., they are fully determined by
\emph{some} combination of the underlying latent parent variables, even
if we do not know their relative contribution\footnote{If \(Y_{tA}\)
  and/or \(Y_{tB}\) contain measurement error, this full measurement
  error is propagated into the composite variable as well.}. Note that
despite \(Y_t = Y_{ta} + Y_{tb}\), \(Y_t\) is not a compositional
variable as defined by Berrie et al.
(\citeproc{ref-berrie2025depicting}{2025}), because the parents
\(Y_{tA}\) and \(Y_{tB}\) are not simultaneously determined. Rather,
\(Y_t\) is a composite derived variable where the parents are
unobservable or latent.\footnote{If you have a strong theoretical reason
  to assume a variable affects one parent but is exogenous to the other,
  you could use it as an instrument to isolate the parent you are
  interested in, e.g., a randomized employment program affecting only
  the second half of the year.} Composite measures present unique
challenges for causal inference (see
\citeproc{ref-berrie2025depicting}{Berrie et al. 2025}), of which we
will highlight two: time-varying confounding and consistency violations.
Time-varying confounding can be a serious problem as Berrie et al.
(\citeproc{ref-berrie2025depicting}{2025, 475}) note:

\begin{quote}
\ldots if the parent variables occur at different moments in time, it is
possible they will have different relationships with the supposed
confounders. Indeed, it is possible that a confounder for one parent
variable may be a mediator for another, leaving no means to
appropriately adjust for confounding without also blocking part of the
true causal effect. Identifying the causal effect of a composite derived
variable, as with any analysis of multiple exposures, therefore requires
no time-varying confounding.
\end{quote}

Consistency refers to the assumption required for making causal
inferences that treatments are well-defined, in particular that only one
unique form of treatment exists at any given measured level of the
exposure variable. This may be violated by composite variables because
one can arrive at the same value of the observed child treatment
variable through multiple combinations of the parent variables; i.e.,
someone who worked half of the past year could have worked only the
first half, only the second half, or some mix of those extremes. This
can produce composite variable bias, because the estimated average
causal effect of the composite variable will be a mean of the average
causal effects of the parent variables that is proportionally weighted
toward the parent(s) with larger variance(s).

A related complication arises when retrospective and instantaneous
measures appear together in the same CLPM, as in our employment--crime
example (past-year employment and past-year offending paired with
measures asked ``currently''). Two variables recorded at the same survey
moment may nonetheless reference different temporal windows, so what
appears as a contemporaneous association can in fact reflect lagged or
asymmetrically-lagged processes that the model is not parameterized to
recover. Where retrospective windows differ across measures (e.g., past
7 days versus past 12 months), the issue compounds: the temporal
ordering within a wave is no longer well-defined, and within-wave
covariances absorb a mixture of true contemporaneous and short-lagged
effects whose decomposition depends on assumptions the data cannot be
used to test (see \citeproc{ref-speyer2025importance}{Speyer et al.
2025} for related issues with concurrent effects in RI-CLPM).

\subsection{Estimands}\label{estimands}

Prior to specifying a statistical model---but after assessing the
limitations of the data---researchers should specify an estimand. For
our first example---the effect of employment on offending which is
threatened by reverse causality---we expect that there is a delay
between when someone is caught offending and when they could be fired
from their job but either employment or offending could have occurred
anywhere throughout the year covered by the measurement. If we decreased
gaps between waves---perhaps to daily or weekly observations---or could
obtain a plausible instrumental variable we could convincingly rule out
reverse causality (\citeproc{ref-jha2026instrumental}{Jha and Nagin
2026}), at least for the contemporaneous effect of employment on
offending, e.g., the ``involvement'' effect in Hirschi's
(\citeproc{ref-hirschi1969causes}{1969}) control theory. With the
available data, however, we can only consider employment effects
occurring over a considerable amount of time, such as the effect of
``commitment'' to conventional lines of action
(\citeproc{ref-hirschi1969causes}{Hirschi 1969}). Accordingly, we
specify our estimand as the effect of proportion of weeks employed in
the prior year (\(X_{t-1}\)) on offending in the following year
(\(Y_t\)). For our second example---the theoretically reciprocal effects
of fear of crime and disorder perceptions---we expect that both
processes of updating could occur quite rapidly---possibly more quickly,
even, than could be captured with \emph{daily} data. Convincingly
determining causal ordering could require at least momentary assessment
(e.g., \citeproc{ref-solymosi2015}{Solymosi, Bowers, and Fujiyama 2015})
or even biophysical measurements in a laboratory, an approach which
would come at a cost of construct and external validity
(\citeproc{ref-esterling2025}{Esterling, Brady, and Schwitzgebel 2025}).
But perhaps we are interested in how this updating process produces
lasting changes in how the environment is perceived and if one side of
the process operates more strongly than the other. Based on this, we
will specify two estimands of interest: the effect of disorder
perceptions measured in the prior year (\(X_{t-1}\)) on fear perceptions
measured in the present year (\(Y_t\)) and vice versa
(\(Y_{t-1} \rightarrow X_t\)). Thus, in both cases, we are confident
about the direction of the lagged and cross-lagged relationships we
expect to observe in our available data but uncertain about the
direction (or even existence) of contemporaneous relationships.

\subsection{Statistical Model}\label{statistical-model}

Once the theoretical model and estimands have been established, one can
proceed to developing a statistical model that accounts for the
constraints from the research design---including ambiguity about causal
directions. Here, path diagrams from the SEM tradition become useful.
SEM diagrams encode both causal relations between variables---indicated
by single-headed arrows pointing from cause to effect---and unanalyzed,
unexplained, or ``nuisance'' associations---indicated by double-headed
arrows pointing between two variables. That is, one explicitly
distinguishes between causal relationships and ambiguous associations.
Variables with any causal path leading into them are termed
``endogenous'' variables and those with no causal path leading into them
are termed ``exogenous''. SEM diagrams depart from most DAGs by
depicting the residuals or unexplained variation of endogenous variables
as latent variables causally affecting the observed variable.

Our recommendation is that the causal part of the SEM diagram should be
recursive, corresponding as closely as possible to the theoretical
model. Reciprocal simultaneous effects should generally be avoided as
they require strong identifying assumptions, e.g., valid and relevant
instrumental variables
(\citeproc{ref-wooldridge2010econometric}{Wooldridge 2010};
\citeproc{ref-jha2026instrumental}{Jha and Nagin 2026}). When the timing
of observation results in ambiguity regarding the causal direction
within any given time period---such as when periods \(t=1.5\) and
\(t=2.5\) were not observed in Figure~\ref{fig-time-dag-unobs}---we
encode this ambiguity as an unanalyzed covariance between the variables
in question. When the timing of observations closely corresponds to
theory, these covariances are unnecessary and the resulting statistical
model similarly corresponds directly to the theoretical model.

For both examples, following these guidelines leads us to the
``classic'' three-wave CLPM depicted in panel (a) of
Figure~\ref{fig-sem-clpm}. The causal relations consist of two sets of
relationships: (a) lagged (autoregressive) effects of variables in one
period on their value in the next period, e.g.,
\(X_{t} \rightarrow X_{t+1}\), and (b) cross-lagged effects of variables
in one period on the other variable in the next period, e.g.,
\(X_{t} \rightarrow Y_{t+1}\). The non-causal relations consist of two
sets of covariances: (a) between the exogenous \(X_1\) and \(Y_1\),
which merely permit the predictors to be correlated as in any typical
regression model, and (b) between the residuals of the endogenous
variables (\(\zeta_t\) and \(\epsilon_t\))\footnote{Omitting the
  residual covariances yields the oft-maligned cross-lagged correlations
  model.}. The residual covariances account for the ambiguous
contemporaneous association between \(X_t\) and \(Y_t\) but cannot tell
us the direction of any causal effect between the two variables.

In principle, the classic CLPM permits identification of the
cross-lagged effects under conventional assumptions common to most
causal inference approaches (see
\citeproc{ref-morgan2014counterfactuals}{Morgan and Winship 2014};
\citeproc{ref-wooldridge2010econometric}{Wooldridge 2010}) and in
particular mediation analysis, such as sequential ignorability
(\citeproc{ref-sobel2008}{Sobel 2008}; \citeproc{ref-imai2010}{Imai,
Keele, and Tingley 2010}). In application, researchers using the classic
CLPM may easily fall prey to a number of common issues that, if
unaddressed, result in misleading inferences. Rather than providing an
exhaustive list of problems that can arise with CLPMs, we focus on the
three key issues we believe are the most important for researchers in
developmental and life course criminology to be aware of due to the
frequency with which they will be encountered and the severity of their
effect on inferences: (a) unobserved heterogeneity, (b) improperly
specified temporal order, and (c) insufficient inter-temporal variation.

\section{Issue 1: Unobserved
heterogeneity}\label{issue-1-unobserved-heterogeneity}

One of the main advantages of drawing on longitudinal data is the
ability to focus exclusively on within-unit variation (i.e., change over
time) and control away any stable, trait-like differences between units
(\citeproc{ref-imai2019should}{Imai and Kim 2019};
\citeproc{ref-allison2017maximum}{Allison, Williams, and Moral-Benito
2017}). For example, by only drawing on within-unit variation, the
family of regression models known as fixed effects models ensure that
any bias emerging from time-constant confounders is accounted for
(although see \citeproc{ref-imai2019should}{Imai and Kim 2019}). In the
traditional cross-lagged panel model, autoregressive paths are supposed
to function similarly to fixed effects, capturing all stable
heterogeneity, thus accounting for time-invariant traits and ensuring
that cross-lagged paths reflect effects on change over time.

However, autoregressive parameters alone are not sufficient to remove
all unobserved stable heterogeneity
(\citeproc{ref-morgan2014counterfactuals}{Morgan and Winship 2014}). Any
stable between-person confounder---that therefore influences both an
outcome \(Y_t\) and a lagged outcome \(Y_{t-1}\)---creates a spurious
association between the two variables that the model incorrectly
attributes to within-unit dynamics. This inflated association then
biases the cross-lagged coefficients
(\citeproc{ref-hamaker2015critique}{Hamaker, Kuiper, and Grasman 2015}),
which no longer reflects \emph{only} effects on change over time. This
implies that the standard CLPM conflates within-unit change with
between-unit differences, and cross-lagged estimates can reflect both
sources of variation.

If autoregressive parameters alone do little to adjust for all
unobserved time-stable confounders, one might contemplate including some
fixed-effects estimator {[}e.g., within-unit demeaning; see Allison
(\citeproc{ref-allison2009fixed}{2009}){]}. However, this strategy
introduces other problems. Imai and Kim
(\citeproc{ref-imai2019should}{2019}) showed that unit fixed effects
only consist of a non-parametric solution to adjusting for all
time-stable heterogeneity under the assumption of no reverse causality
or causal dynamics---that is, that the past outcome does not affect
either the current treatment or outcome. Such dynamic relationships are
precisely the issue that CLPMs attempt to address, so they cannot be
assumed away. In addition, including fixed effects alongside
autoregressive parameters in a dynamic panel model can lead to the
well-known \emph{Nickell bias} (\citeproc{ref-nickell1981biases}{Nickell
1981}): estimates are downwardly biased with few periods (i.e., small
\(T\)) because of the connection between the lagged outcome and the
fixed effects (\citeproc{ref-wooldridge2010econometric}{Wooldridge
2010}). In the econometrics literature, this is usually solved with the
Arellano-Bond estimator (\citeproc{ref-arellano1991some}{Arellano and
Bond 1991}), which removes Nickell bias and adjusts for unobserved
stable heterogeneity---but it yields biased estimates with small samples
and when the autoregressive parameter is near 1 (i.e., the outcome is
highly stable over time).

Stable unobserved heterogeneity is therefore an important issue. The
standard CLPM relies on autoregressive parameters to capture temporal
stability, but this does not properly ensure that only within-unit
variation is taken into account. Any unmodeled time-stable confounder
can bias the cross-lagged parameters, which end up conflating
within-unit change with between-unit differences. The Arellano-Bond
estimator of the dynamic panel model can solve this issue by including
fixed effects, but it has other limitations.

\subsection{Solutions}\label{solutions}

While the traditional CLPM cannot properly account for unobserved stable
heterogeneity, some relatively recent extensions of the CLPM yield
robust estimators that draw only on within-unit variation and
effectively control for all unobserved time-stable traits
(\citeproc{ref-leszczensky2019deal}{Leszczensky and Wolbring 2019}).
Specifically, Allison et al.'s (\citeproc{ref-allison2017maximum}{2017})
maximum likelihood structural equation fixed-effects dynamic panel
method (ML-SEM) and Hamaker et al.'s
(\citeproc{ref-hamaker2015critique}{2015}) random intercept cross-lagged
panel model (RI-CLPM) model reciprocality and account for unobserved
time-stable features that impact both \(Y_t\) and \(X_t\). It is our
position that the ML-SEM method is better suited to cases of reverse
causality (e.g., our work and crime example) and RI-CLPM is better
suited to cases of theoretical reciprocality (e.g., our disorder and
fear example).

\subsubsection{Fixed-effects dynamic panel model with maximum
likelihood}\label{fixed-effects-dynamic-panel-model-with-maximum-likelihood}

In ML‐SEM (\citeproc{ref-allison2017maximum}{Allison, Williams, and
Moral-Benito 2017}), a dynamic panel model (i.e., with a lagged
dependent variable) is estimated as a structural equation model that
partials out all unobserved time-stable heterogeneity. This is
accomplished using a latent random intercept factor (often denoted
\(\alpha\)) with a factor loading constrained to 1 on every observed
outcome \(Y_t\). This latent random intercept factor is the same as the
one often included in latent growth curve models (e.g., a random
intercept reflecting the starting point of latent trajectories) and
represents all between-unit differences. By fixing the factor loadings
of the random intercept to 1, the model effectively ``demeans'' each
unit's scores and thus conditions out all time‐constant heterogeneity.
Crucially, this latent random intercept factor is allowed to freely
covary with all the time‐varying exogenous predictors (including any
predetermined regressors and the initial outcome \(Y_1\)).\footnote{This
  is a correlated random‐effects model or latent Mundlak model
  (\citeproc{ref-mundlak1978pooling}{Mundlak 1978}): by letting
  \(\alpha\) correlate without restriction with all \(X\), the SEM
  mimics a conventional fixed‐effects specification. In contrast, a
  conventional uncorrelated random‐effects model would force \(\alpha\)
  to be independent of \(X\).} As Allison, Williams, and Moral-Benito
(\citeproc{ref-allison2017maximum}{2017}) note, allowing \(\alpha\) to
covary with \(X\)'s mimics fixed-effects model that account for all
unobserved confounders that are time-stable within the observation
period. In contrast with the conventional CLPM, ML-SEM focuses on
estimating one side of the cross-lagged relationship (e.g., cross-lagged
and/or contemporaneous effects of \(X\) on \(Y\)) while accounting for
the opposite reciprocal effect (the lagged effect of \(Y\) on \(X\))
using a covariance between past outcome (\(Y_{t-1}\)) and current
predictor (\(X_t\)).\footnote{If cross-lagged or contemporaneous effects
  of the outcome \(Y_t\) (e.g., \(Y_{t-1}\) or \(Y_t\)) on the predictor
  \(X_t\) are also of substantive interest, one needs to fit a separate
  ML-SEM with that specification---taking care that these are
  theoretically and empirically coherent, e.g., that both do not include
  contemporaneous effects. If this is the case, however, we advise
  instead fitting an RI-CLPM described below that is better suited to
  studying cases of theoretical reciprocality.} This makes it most
applicable to cases of reverse causality, where the reciprocal feedback
effect is not of substantive interest, such as in our work and crime
example.

In implementation, the entire system of autoregressive and cross‐lagged
equations is estimated as one SEM. For example, Figure~\ref{fig-mlsem}
illustrates an ML-SEM with four waves with both contemporaneous and
lagged effects of \(X\) on \(Y\), that is each outcome (\(Y_t\)) is
regressed on its lagged value (\(Y_{t-1}\)), on a contemporaneous
(\(X_t\)) and lagged value of the predictor (\(X_{t-1}\)).\footnote{The
  ML-SEM with lagged effects requires at least four waves of data to be
  implemented.} \(\alpha\) represents the latent random intercept factor
reflected by all observations \(Y_t\) (except for the initial
observation \(Y_{t=1}\)) with factor loadings constrained to 1, and all
\(X_t\) are allowed to correlate with \(\alpha\) (and with each other
and with \(Y_{t=1}\)). In this setup, the first observation \(Y_{t=1}\)
is treated like any exogenous covariate. Importantly, observed
time-varying regressors (\(X_t\)) are specified as predetermined---that
is, they are allowed to correlate with past residuals (i.e., with
\(\epsilon_{t-1}\))---reflecting that predictors may be influenced by
prior outcomes. The presence of contemporaneous effects means a strict
assumption of no contemporaneous reverse causality (i.e.,
\(Y_t \rightarrow X_t\) does not exist); if reverse causality is a
risk---such as in our work and crime example---the researcher would
replace these contemporaneous effects with covariances to reflect the
causal ambiguity. The resulting model is asymptotically equivalent to
the Arellano-Bond method (\citeproc{ref-arellano1991some}{Arellano and
Bond 1991}; \citeproc{ref-allison2017maximum}{Allison, Williams, and
Moral-Benito 2017})\footnote{It however requires no stationarity
  assumption on initial conditions given that initial observations
  \(Y_{t=1}\) and \(X_{t=1}\) are treated as additional exogenous
  predictors and does not suffer from issues regarding instrument
  proliferation (\citeproc{ref-allison2017maximum}{Allison, Williams,
  and Moral-Benito 2017}).}, but can additionally handle missing data
via full information maximum likelihood, include measurement models for
multiple indicators, incorporate time-constant covariates (uncorrelated
with the fixed effects), and fit period-specific coefficients on
regressors to account for practical issues like uneven observation
periods or theoretical concerns like effects that differ by life stage.

A number of recent studies used ML-SEM to examine (potentially)
reciprocal relationships. In life course and developmental criminology,
researchers have examined cross-lagged relationships between bullying
victimization and perpetration (\citeproc{ref-kim2023}{Kim et al.
2023}), delinquency and arrest (\citeproc{ref-demarco2025}{DeMarco,
Leppard, and Lindsay 2025}), and offending and peer delinquency
(\citeproc{ref-unnever2020}{Unnever and Chouhy 2020}). In other areas of
criminology, ML-SEM has been applied to examine de-policing and homicide
rates in American cities (\citeproc{ref-rosenfeld2019}{Rosenfeld and
Wallman 2019}), family conflict and violence following release from
prison (\citeproc{ref-stansfield2020}{Stansfield et al. 2020})---and by
the present authors to Airbnb short term lettings and crime in London
neighborhoods (\citeproc{ref-lanfear2024promise}{Lanfear and Kirk 2024})
and legal cynicism, police legitimacy, and intrusive policing in Brazil
(\citeproc{ref-oliveira2025}{Oliveira 2025}).

\begin{figure}

\centering{

\pandocbounded{\includegraphics[keepaspectratio]{./jdlcc_draft_deidentified_files/figure-pdf/fig-mlsem-1.pdf}}

}

\caption{\label{fig-mlsem}Maximum likelihood structural equation
fixed-effects dynamic panel model with contemporaneous and lagged
effects}

\end{figure}%

\subsubsection{Random intercept cross-lagged panel
model}\label{random-intercept-cross-lagged-panel-model}

Another alternative to the traditional cross-lagged panel model that
controls for unobserved stable heterogeneity is Hamaker et al.'s
(\citeproc{ref-hamaker2015critique}{2015}) RI-CLPM. To ensure that the
model only takes into account change over time, the RI-CLPM models
within-unit and between-unit variances as separate attributes. Similar
to Allison et al.'s (\citeproc{ref-allison2017maximum}{2017}) ML-SEM,
RI-CLPM as depicted in Figure~\ref{fig-riclpm} estimates a latent
variable that captures all stable variation between units---again
drawing on the latent growth curve model specification, this latent
random intercept factor is defined by all observations \(Y_t\) over time
with factor loadings constrained to 1, essentially creating a random
intercept for \(Y\) that varies between units. While the ML-SEM allows
the latent random intercept factor to covary with predetermined
regressors---which implies treating it as a fixed effect that captures
all unobserved stable heterogeneity---the RI-CLPM instead repeats the
process and creates another latent random intercept factor defined by
all observations \(X_t\) over time that varies between units. The two
latent random intercept factors for \(Y_t\) and \(X_t\) (represented by
\(\alpha_Y\) and \(\alpha_X\) in Figure~\ref{fig-riclpm}) are allowed to
covary, which ensures that they capture all unobserved stable
heterogeneity across both variables.

\begin{figure}

\centering{

\pandocbounded{\includegraphics[keepaspectratio]{./jdlcc_draft_deidentified_files/figure-pdf/fig-riclpm-1.pdf}}

}

\caption{\label{fig-riclpm}Random Intercept Cross-Lagged Panel Model
(Hamaker et al., 2015)}

\end{figure}%

To ensure that only within-unit variation is taken into account in the
structural model, the RI-CLPM relies on another set of latent variables:
For each variable \(Y_t\) and \(X_t\), a new latent variable is
estimated with a factor loading constrained to 1. These latent variables
are represented in the diagram in Figure~\ref{fig-riclpm} as \(W_{Y_t}\)
and \(W_{X_t}\). This implies that every observed variable \(Y_t\) and
\(X_t\) is determined by two latent variables: one capturing all
time-stable differences between units (i.e., \(\alpha_Y\) and
\(\alpha_X\)) and one capturing everything else, i.e., within-unit
differences over time (i.e., \(W_{Y_t}\) and \(W_{X_t}\)).
Autoregressive and cross-lagged parameters are then included linking
\(W_{Y_t}\) and \(W_{X_t}\)---resembling the traditional cross-lagged
panel model, but exclusively focused on within-unit variation over time.
By allowing both latent random intercept factors to covary, the RI-CLPM
ensures that all unobserved stable heterogeneity is accounted for; and
by estimating autoregressive and cross-lagged effects on latent
variables capturing only within-unit variation, the model ensures that
cross-lagged estimates reflect only change over time. In contrast to
ML-SEM, which focuses on one side of the reciprocal relationship,
RI-CLPM produces a symmetrical model that considers both reciprocal
pathways at once and is thus better suited to studying theoretical
reciprocality---such as our example of the reciprocal relationship
between disorder and fear. Like ML-SEM, the RI-CLPM can also be extended
to other common features in SEMs, such as multiple indicators, covariate
adjustment, period-specific effects, and multiple-group analysis
(\citeproc{ref-mulder2021three}{Mulder and Hamaker 2021}) as well as
simultaneous reciprocal effects under strong identifying assumptions
(\citeproc{ref-speyer2025importance}{Speyer et al. 2025};
\citeproc{ref-muthuxe9n2024}{Muthén and Asparouhov 2024}).

The RI-CLPM has already seen wide use in developmental and life course
criminology, likely due in part to its proximity to psychology where the
method originates. In the \emph{Journal of Developmental and Life Course
Criminology} alone, it has been used to examine (a) reciprocality
between psychosocial maturity, resistance to peer influence, and
antisocial influence by peers (\citeproc{ref-forney2019identity}{Forney
and Ward 2019}), (b) time spent in correctional confinement and
perceptions of legitimacy and legal cynicism
(\citeproc{ref-chen2026}{Chen and Nichols 2026}), and (c) mental health
and self-reported offending (\citeproc{ref-wiesner2023}{Wiesner et al.
2023}). The RI-CLPM appears less commonly used outside of
psychologically-oriented criminology.

\section{Issue 2: Temporal Order}\label{issue-2-temporal-order}

Temporal order presents two key issues for studying reciprocal effects.
The first is that researchers may fail to consider time in their
theoretical model to begin with. This may result in choosing an
estimator before specifying---or perhaps without even knowing---their
estimand. If yearly panel data are available, they simply estimate the
effect of some past-year measure on a present year measure. This may
also result in researchers assuming that if they fail to detect an
effect in one lag structure it implies effects do not exist at other
temporal lags. Researchers may instead assume there exists one ``true
causal effect'' between two variables which is ambiguous in timing and
set out to use a model to determine at what timing this effect exists.
This raises the question of whether different causal timings could
possibly correspond to comparable or even plausible processes of
interest, and secondarily whether any model is able to identify multiple
causal timings simultaneously.

Improper specifications of temporal order may have serious consequences
for estimation. Simulations specifying a causal effect that exists at
only one particular time lag reveal that incorrect specifications of
temporal lags in statistical models with fixed effects can produce
estimates with severe bias including flipped signs. As Vaisey and Miles
(\citeproc{ref-vaisey2017you}{2017}) document, if Equation \ref{eq1} is
the true reduced form model of interest (\(x\) has only a
contemporaneous effect) and one instead specifies Equation \ref{eq2} in
their statistical model (\(x\) with only a lagged effect), then
\(E(\beta^*) = -0.5\beta\), i.e., ``a positive contemporaneous effect of
\(x\) gets transformed into an artifactually negative lagged effect''
(\citeproc{ref-allison2022getting}{Allison 2022}).

\begin{align}
y_{ti} &= \beta x_t + \alpha_i + e_{it} \label{eq1}\\
y_{ti} &= \beta^* x_{t-1} + \alpha_i + e_{it} \label{eq2}
\end{align}

This is possibly the least appreciated and most serious issues facing
CLPMs as it may produce misleading results that are interpreted as
evidence not just against the theoretical proposition being tested but
even for its opposite; to use one of our examples, this could result in
research finding past employment increases future offending even when it
has only a preventative contemporaneous effect. Note that this is not
the result of signs differing between estimands but rather a biased
estimate resulting from misspecification in equation 2, specifically the
assumption of no contemporaneous effects between \(y_t\) and \(x_t\). In
our example, this is would be assuming crime has no effect on employment
within the same year, i.e., assuming away reverse causality. Research
that finds unexpectedly negative estimates from CLPMs when theory
suggests positive effects may be the result of this bias (e.g.,
\citeproc{ref-lanfear2020broken}{Lanfear, Matsueda, and Beach 2020,
114}; \citeproc{ref-allison2022getting}{Allison 2022}).

Applied to our employment--crime example: if Hirschi's
(\citeproc{ref-hirschi1969causes}{1969}) involvement hypothesis is
correct and employment reduces offending contemporaneously by occupying
time, but a researcher models employment only with a one-year lag, the
resulting estimate need not just be attenuated; it can flip positive,
suggesting employment increases future offending. A published
criminological example of this dynamic comes from Ousey, Wilcox, and
Fisher (\citeproc{ref-ousey2011something}{2011}). They used a
fixed-effects SEM to estimate cross-lagged effects between victimization
and offending, and reported negative cross-lagged coefficients in both
directions, contradicting a large body of research that had consistently
found a positive association between victimization and offending. As
Allison (\citeproc{ref-allison2022getting}{2022}) later observed, these
negative lagged estimates are precisely what Vaisey and Miles
(\citeproc{ref-vaisey2017you}{2017}) had warned about: when the true
effects are contemporaneous (or operate at a lag substantially shorter
than observed measures) and the model includes a large (annual in this
case) lag, a flip sign occurs. It is plausible that victimization and
offending are actually positively related, and the negative coefficients
reported by Ousey, Wilcox, and Fisher
(\citeproc{ref-ousey2011something}{2011}) are a flip sign due to an
inappropriate lag or covariance specification. This study is otherwise
methodologically thoughtful, as authors applied state-of-the-start
methods at the time, when the sign flip issue was not as well
recognized.

It is difficult to know how widespread this issue is in the
criminological literature---in part because model specifications are not
always clearly described---but we believe the potential consequences
warrant a systematic examination of its frequency and consequences given
that lag-only specifications are a common default (e.g.,
\citeproc{ref-allison2017maximum}{Allison, Williams, and Moral-Benito
2017}) and that research may be more likely to be published when it
produces novel findings such as those resulting from unexpected sign
flips.

\subsection{Solutions}\label{solutions-1}

As we have stated earlier, there is no ``one true causal effect'' but
rather a number of causal effects corresponding to different estimands
researchers choose. In this regard, what is important is specifying the
correct lag structure \emph{for the chosen estimand}. A null effect for
a particular time lag may be compatible with a non-null effect at a
different time period; a study that finds no effect of an additional
pint of beer on decision-making with a lag of one day provides little
information on the effect at one hour. It is unlikely a researcher has
no theoretical expectations about what causal timings are of interest.
If this is the case, however, the researcher is not prepared to engage
in a causal analysis. Determining which lags are of theoretical interest
is a question for literature reviews and exploratory analyses.

The prior equations specify different estimands of interest, though both
make the strong assumption there is no reverse causality between \(x_t\)
and \(y_t\). If the estimand is the contemporaneous effect of \(x_{t}\)
on \(y_t\) (Equation \ref{eq1}) and there is concern about reverse
causality, some auxiliary information is necessary such as an
instrumental variable for \(x_t\). Without a theoretical justification
otherwise, it is reasonable to default to including a cross-lagged
effect as well (\(x_{t-1} \rightarrow y_t\)). If also accounting for
unobserved heterogeneity, this necessitates use of the models above
which address Nickell bias occurring in dynamic panel models with
unobserved effects. If the estimand of interest is the lagged effect of
\(x_{t-1}\) on \(y_t\) (Equation \ref{eq2}) however, then reverse
causality could be addressed in a cross-lagged model with a covariance
between each \(x_t\) and \(y_t\) to address this nuisance reciprocality
that may bias the estimate of interest
(\citeproc{ref-allison2022getting}{Allison 2022}). This is a more robust
specification if the parameter of interest is \(\beta x_{t-1}\) above,
but it cannot estimate \(\beta x_{t}\) or adjudicate between the
different directions of contemporaneous effects. However, provided
strong theoretical expectations that reverse causality is absent or
minimal, researchers may fit models with contemporaneous effects.

For example, Speyer et al. (\citeproc{ref-speyer2025importance}{2025})
used measurements corresponding to different within-period timeframes
for victimization (past year) and negative mental states (past week) to
justify contemporaneous effects from the longer to shorter
timeframe---i.e., theory regarding measurement permits establishing
within-wave temporal order. Similarly, Lanfear and Kirk
(\citeproc{ref-lanfear2024promise}{2024}) justified ML-SEM models with
both lagged and contemporaneous effects of short-term Airbnb lettings on
crime under the assumption that lettings present immediate criminal
opportunities but the letting market adjusts slowly to local changes in
crime---i.e., assumptions about causal order came from theory.

\section{Issue 3: Inter-temporal
variation}\label{issue-3-inter-temporal-variation}

Many of the focal causes and effects in developmental and life course
criminology are relatively stable over time in some sense, such as
psychological constructs (e.g., self-control) and life-course statuses
(e.g., marriage). By definition, in cross-lagged panel models the
outcomes serve as predictors of later outcomes---both as lags (e.g.,
\(Y_{t-1} \rightarrow Y_t\)) and cross-lags (e.g.,
\(Y_{t-1} \rightarrow X_t\)). If an outcome is relatively stable over
time---that is, strongly dependent on its past values---there may be
little remaining to be explained by the predictors of interest after
accounting for its prior value. Put another way, as the correlation
between past and present values increases
(\(\rho(Y_t,Y_{t-1}) \rightarrow 1\)) the remaining conditional variance
that can be explained by other covariates shrinks
(\(Var(Y_t|Y_{t-1}) \rightarrow 0\)).\footnote{Note that if the
  variables are measured with error, what is relevant is the underlying
  correlation between the error-free latent outcomes; when measurement
  error is present, it becomes an increasingly large proportion of the
  conditional variance as the true correlation increases. This may
  result in even more imprecise estimates while also attenuating
  estimates as the outcomes also serve as regressors. This may be
  addressed using measurement models or errors-in-variables estimators.}
This is an imperfect multicollinearity problem: estimates may remain
unbiased and consistent but imprecise and even unstable or empirically
underidentified when the issue is extreme. This is exacerbated in
estimators with unobserved effects as relatively stable outcomes may be
largely explained by the time-stable fixed effects or correlated random
intercepts.\footnote{One may also include time-invariant covariates in
  unobserved effects models using a hybrid random effects approach,
  however the estimates will correspond to between-individual effects
  which may differ from effects of within-individual changes and must be
  assumed to be uncorrelated with unobserved effects.}

Many of the effects of central interest in developmental and life-course
criminology are theorized to accumulate gradually rather than switch on
at a moment in time, which makes them difficult to study over short
observation windows. Consider the employment effect in our running
example. The near-immediate ``involvement'' effect of employment
occupying one's time is distinct from the slower ``commitment'' effect
(\citeproc{ref-hirschi1969causes}{Hirschi 1969}), in which the
accumulation of conventional ties and responsibilities gradually
strengthens social control and thereby reduces offending
(\citeproc{ref-sampson1995crime}{Sampson and Laub 1995}). It is no
accident that the foundational within-person evidence for such
turning-point effects was drawn from data following individuals across
decades (e.g., \citeproc{ref-sampson2006does}{Sampson, Laub, and Wimer
2006}): the cumulative change attributable to a slowly-building bond is
difficult to detect over a short span because too little of it has yet
accrued. These issues are common in developmental studies of individuals
but also commonly found in other criminological literatures using panel
data (see \citeproc{ref-lanfear2020broken}{Lanfear, Matsueda, and Beach
2020}).

Studies with finer temporal resolution make clear just how extended
these processes are. Using Norwegian register data, Skardhamar and
Savolainen (\citeproc{ref-skardhamar2014changes}{2014}) found that the
decline in offending around job entry largely precedes employment rather
than following it, and Skardhamar, Monsbakken, and Lyngstad
(\citeproc{ref-skardhamar2014crime}{2014}) report that desistance tends
to begin years before marriage. Whatever one concludes from these
anticipation patterns, they demonstrate that the relevant within-person
change, in this case, is spread across a long-term trajectory. A
researcher who observes such a process with only two or three annual
waves sees a relatively stable outcome over that span, and
correspondingly little within-person variation for the effect of
interest to explain---not because there are no effects, but because the
observed temporal variation is too short relative to the timescale on
which the process unfolds. For example, Forney and Ward
(\citeproc{ref-forney2019identity}{2019}) fit a RI-CLPM with annual lags
to even waves of data from the Pathways to Desistance study and reported
significant autoregressive parameters across several constructs (e.g.,
identity, peer influences, delinquency). This could indicate (and it
likely does) that these constructs are highly stable over time; but
equally, this could also indicate low within-unit variation captured by
annual lags, and other temporal aggregations could yield different
results.

\subsection{Solutions}\label{solutions-2}

Ideally, researchers should account for the possibility of low
inter-temporal variation during design, but in general choice of
temporal resolution should be guided by strong theory
(\citeproc{ref-hamaker2026}{Hamaker 2026};
\citeproc{ref-nivette2026}{Nivette and Eisner 2026}). Researchers
interested in slow processes and relatively stable outcomes can
obviously opt to collect data over long time periods. If doing so is
cost prohibitive or could prevent answering a timely research question,
there are at least two other options. First, if the normally time-stable
factor is affected by external factors---including those that may be
manipulated by researchers---then researchers can collect data on
individuals exposed to shocks in these factors---or embed an
intervention in their study to manipulate the external factor. This then
opens the possibility of using an instrumental variables estimator to
estimate a local average treatment effect
(\citeproc{ref-jha2026instrumental}{Jha and Nagin 2026}). Second, even
if the time stable factor is not responsive to exogenous shocks, the
factor may be less stable among some subpopulations. In this case,
researchers might oversample the more rapidly changing subpopulation to
maximize variation, while remaining aware that the effect identified in
that group may differ from that in the larger stable population.

Researchers using already-collected data may consider a number of
approaches. First, they may consider estimands corresponding to shorter
time aggregations or longer time intervals. Disaggregation may help as
low variation at a larger temporal unit (e.g., a year) may mask higher
variation in smaller units (e.g., a month). However, unless such an
estimator is intractable with many time periods, researchers should
already default to disaggregated analysis as longer-term causal effects
can always be captured using models with multiple lags of shorter
periods. In extreme situations, researchers may consider skipping waves
if it makes sense theoretically (e.g., skipping a mid-year wave when
considering an annual process) or they are facing serious empirical
underidentification problems related to a particular wave (e.g., a very
short observation interval resulted in no change between waves). In any
case, researchers are again advised not to substitute a good question
that is difficult to answer (i.e., a theoretically-motivated estimand)
for a bad question that is easy to answer (i.e., anything else).

\section{Conclusion}\label{conclusion}

Our key recommendations for researchers studying potentially reciprocal
relationships with repeated observations may be summarized in three
points.

First, it is important to separate theory problems from estimation
problems. A clear theoretical model should be constructed first based on
deep knowledge of the substantive research area and the data generation
process. From this, researchers should then determine what, if any,
theoretically-relevant estimand can plausibly be estimated using the
available data. With a theory and estimand in hand, researchers should
then choose an appropriate estimator. Readers may also apply this logic
when reading research, asking if the stated or implied estimand is one
of substantive interest or one that corresponds to a plausible
mechanism. If not, a strong estimation strategy cannot make up for a
theoretical deficit.

Second, researchers should think carefully about whether the data they
have are suitable for answering their questions to begin with---that is,
they should be prepared to stop after the first step above. Conventional
longitudinal panel survey data are better equipped to test slow and
lasting effects proposed by developmental and life course theories
(e.g., trajectories and turning points) than rapid or transient
processes from cognitive and interactional theories (e.g.,
decision-making and group processes). For testing theories that connect
long-term and short-term processes, researchers should consider
incorporating other forms of data collection into panel studies---e.g.,
experiments (\citeproc{ref-farrington2009longitudinal}{Farrington,
Loeber, and Welsh 2009}), experience sampling or momentary assessments
(\citeproc{ref-solymosi2015}{Solymosi, Bowers, and Fujiyama 2015};
\citeproc{ref-browning2024leveraging}{Browning et al. 2024}),
experiential cascade analysis (\citeproc{ref-roman2026}{Roman 2026}), or
retrospective calendar methods (\citeproc{ref-apel2017}{Apel and Horney
2017}; \citeproc{ref-wikstruxf6m2024}{Wikström, Treiber, and Roman
2024}). Panel data with narrow observation intervals are also poorly
suited for testing slow processes; high intertemporal correlations and
proportions of variance explained may be signals there is insufficient
change over time to produce precise estimates. Researchers and readers
should also consider if they are using composite period variables and
how this may impact the validity of their designs and estimates; these
variables may correspond to unusual estimands, violate consistency
assumptions for making causal inferences, and make it difficult to
plausibly address reverse causality.

Third, researchers should default to robust estimators. If the estimand
is a lagged effect, include a contemporaneous covariance to adjust for
potential reverse causality. If the treatment and outcome of interest
vary (sufficiently) over time, adjust for unobserved heterogeneity using
ML-SEM or RI-CLPM. If, on the one hand, one has clearly defined
dependent and independent variables, but reverse causality is an
important nuisance that needs to be handled, the ML-SEM is the better
choice of estimator; if, on the other hand, the substantive question is
one of theoretical reciprocality, then the RI-CLPM is the preferred
estimator. If the estimand is a contemporaneous effect (or the effect of
a stable trait) or there is insufficient variation to use these robust
estimators, researchers should state outright the strong assumptions
they are making---which, importantly, may still be the best that can be
done to answer an important question until better data and designs are
available. When evaluating research, readers should ask whether the
chosen estimator is consistent with the theoretical model and robust to
common issues we describe. In particular, estimates opposite in sign of
theoretical expectations should invite skepticism, scrutiny, and---in
particular---replication; a literature search for studies with fixed
effects, lagged dependent variables, and no contemporaneous covariance
would produce a reservoir of studies to replicate with robust
estimators.

Finally, while this work focuses on issues specific to studies positing
reciprocal effects, researchers and readers should also be attentive to
key issues that apply to all quantitative studies testing theories such
as plausibly establishing (conditional) ignorability, consistency, and
absence of interference (\citeproc{ref-morgan2014counterfactuals}{Morgan
and Winship 2014}; \citeproc{ref-rubin1986comment}{Rubin 1986}). In
longitudinal contexts, this includes considering dynamic effects (i.e.,
lagged outcomes) and time-varying confounding. Where these are difficult
to address, exogenous shocks, field interventions, and embedded lab
experiments may provide plausible identification strategies. General
training in principles for causal inference thus serves as a foundation
for strong longitudinal research designs. Fortunately, barriers to entry
for developing these skills have rapidly eroded with the proliferation
of accessible introductory materials (e.g.,
\citeproc{ref-morgan2014counterfactuals}{Morgan and Winship 2014};
\citeproc{ref-huntington2021effect}{Huntington-Klein 2021};
\citeproc{ref-cunningham2021causal}{Cunningham 2021};
\citeproc{ref-llaudet2023data}{Llaudet and Imai 2023}).

\section*{References}\label{references}
\addcontentsline{toc}{section}{References}

\phantomsection\label{refs}
\begin{CSLReferences}{1}{0}
\bibitem[\citeproctext]{ref-allison2009fixed}
Allison, Paul D. 2009. \emph{Fixed Effects Regression Models}. SAGE
publications.

\bibitem[\citeproctext]{ref-allison2022getting}
---------. 2022. {``Getting the Lags Right -- a New Solution.''}
\emph{Statistical Horizons}.
\url{https://statisticalhorizons.com/getting-the-lags-right-a-new-solution/.}

\bibitem[\citeproctext]{ref-allison2017maximum}
Allison, Paul D, Richard Williams, and Enrique Moral-Benito. 2017.
{``Maximum Likelihood for Cross-Lagged Panel Models with Fixed
Effects.''} \emph{Socius} 3.

\bibitem[\citeproctext]{ref-apel2017}
Apel, Robert, and Julie Horney. 2017. {``How and Why Does Work Matter?
Employment Conditions, Routine Activities, and Crime Among Adult Male
Offenders.''} \emph{Criminology} 55 (2): 307--43.
\url{https://doi.org/10.1111/1745-9125.12134}.

\bibitem[\citeproctext]{ref-arellano1991some}
Arellano, Manuel, and Stephen Bond. 1991. {``Some Tests of Specification
for Panel Data: Monte Carlo Evidence and an Application to Employment
Equations.''} \emph{The Review of Economic Studies} 58 (2): 277--97.

\bibitem[\citeproctext]{ref-becker1963outsiders}
Becker, Howard S. 1963. \emph{Outsiders: Studies in the Sociology of
Deviance}.

\bibitem[\citeproctext]{ref-berrie2025depicting}
Berrie, Laurie, Kellyn F Arnold, Georgia D Tomova, Mark S Gilthorpe, and
Peter WG Tennant. 2025. {``Depicting Deterministic Variables Within
Directed Acyclic Graphs: An Aid for Identifying and Interpreting Causal
Effects Involving Derived Variables and Compositional Data.''}
\emph{American Journal of Epidemiology} 194 (2): 469--79.

\bibitem[\citeproctext]{ref-bollen2013eight}
Bollen, Kenneth A, and Judea Pearl. 2013. {``Eight Myths about Causality
and Structural Equation Models.''} In \emph{Handbook of Causal Analysis
for Social Research}, 301--28. Springer.

\bibitem[\citeproctext]{ref-browning2024leveraging}
Browning, Christopher R, Nicolo P Pinchak, Catherine A Calder, and
Bethany Boettner. 2024. {``Leveraging Experience Sampling/Ecological
Momentary Assessment for Sociological Investigations of Everyday
Life.''} \emph{Annual Review of Sociology} 50 (1): 41--59.

\bibitem[\citeproctext]{ref-brunton2011untangling}
Brunton-Smith, Ian. 2011. {``Untangling the Relationship Between Fear of
Crime and Perceptions of Disorder: Evidence from a Longitudinal Study of
Young People in {E}ngland and {W}ales.''} \emph{The British Journal of
Criminology} 51 (6): 885--99.

\bibitem[\citeproctext]{ref-chalfin2011}
Chalfin, Aaron, and Steven Raphael. 2011. {``Work and Crime.''} In,
edited by Michael Tonry, 444--76. Oxford University Press.
\url{https://doi.org/10.1093/oxfordhb/9780195395082.013.0015}.

\bibitem[\citeproctext]{ref-chen2026}
Chen, Frances R., and Cassidey Nichols. 2026. {``Legal Attitudes in
Transition: Bidirectional Links Between Correctional Time and the
Development of Perceived Legitimacy and Legal Cynicism.''} \emph{Journal
of Developmental and Life-Course Criminology} 12 (1): 10.
\url{https://doi.org/10.1007/s40865-026-00298-9}.

\bibitem[\citeproctext]{ref-crutchfield2014get}
Crutchfield, Robert D. 2014. \emph{Get a Job: Labor Markets, Economic
Opportunity, and Crime}. New York University Press.

\bibitem[\citeproctext]{ref-cunningham2021causal}
Cunningham, Scott. 2021. \emph{Causal Inference: The Mixtape}. Yale
university press.

\bibitem[\citeproctext]{ref-demarco2025}
DeMarco, Laura M., Tom R. Leppard, and Sadé L. Lindsay. 2025.
{``Intersectional Bonds: Delinquency, Arrest, and Changing Family Social
Capital During Adolescence.''} \emph{Journal of Marriage and Family} 87
(2): 505--26. \url{https://doi.org/10.1111/jomf.13029}.

\bibitem[\citeproctext]{ref-driver2017continuous}
Driver, Charles C, Johan HL Oud, and Manuel C Voelkle. 2017.
{``Continuous Time Structural Equation Modeling with r Package Ctsem.''}
\emph{Journal of Statistical Software} 77: 1--35.

\bibitem[\citeproctext]{ref-esterling2025}
Esterling, Kevin M., David Brady, and Eric Schwitzgebel. 2025. {``The
Necessity of Construct and External Validity for Deductive Causal
Inference.''} \emph{Journal of Causal Inference} 13 (1).
\url{https://doi.org/10.1515/jci-2024-0002}.

\bibitem[\citeproctext]{ref-fagan1999}
Fagan, Jeffrey, and Richard B. Freeman. 1999. {``Crime and Work.''}
\emph{Crime and Justice} 25 (January): 225--90.
\url{https://doi.org/10.1086/449290}.

\bibitem[\citeproctext]{ref-farrington2009longitudinal}
Farrington, David P, Rolf Loeber, and Brandon C Welsh. 2009.
{``Longitudinal-Experimental Studies.''} \emph{Handbook of Quantitative
Criminology}, 503--18.

\bibitem[\citeproctext]{ref-forney2019identity}
Forney, Megan, and Jeffrey T Ward. 2019. {``Identity, Peer Resistance,
and Antisocial Influence: Modeling Direct and Indirect Causes of
Desistance.''} \emph{Journal of Developmental and Life-Course
Criminology} 5 (1): 107--35.

\bibitem[\citeproctext]{ref-giordano2002}
Giordano, Peggy C., Stephen A. Cernkovich, and Jennifer L. Rudolph.
2002. {``Gender, Crime, and Desistance: Toward a Theory of Cognitive
Transformation.''} \emph{American Journal of Sociology} 107 (4):
990--1064. \url{https://doi.org/10.1086/343191}.

\bibitem[\citeproctext]{ref-gollob1987taking}
Gollob, Harry F, and Charles S Reichardt. 1987. {``Taking Account of
Time Lags in Causal Models.''} \emph{Child Development} 58 (1): 80--92.

\bibitem[\citeproctext]{ref-hamaker2026}
Hamaker, Ellen L. 2026. {``The Within-Between Dispute in Cross-Lagged
Panel Research and How to Move Forward.''} \emph{Psychological Methods}
31 (1): 56--76. \url{https://doi.org/10.1037/met0000600}.

\bibitem[\citeproctext]{ref-hamaker2015critique}
Hamaker, Ellen L, Rebecca M Kuiper, and Raoul PPP Grasman. 2015. {``A
Critique of the Cross-Lagged Panel Model.''} \emph{Psychological
Methods} 20 (1): 102.

\bibitem[\citeproctext]{ref-hirschi1969causes}
Hirschi, Travis. 1969. \emph{Causes of Delinquency}. UC Berkeley Press.

\bibitem[\citeproctext]{ref-huntington2021effect}
Huntington-Klein, Nick. 2021. \emph{The Effect: An Introduction to
Research Design and Causality}. Chapman; Hall/CRC.

\bibitem[\citeproctext]{ref-imai2010}
Imai, Kosuke, Luke Keele, and Dustin Tingley. 2010. {``A General
Approach to Causal Mediation Analysis.''} \emph{Psychological Methods}
15 (4): 309--34. \url{https://doi.org/10.1037/a0020761}.

\bibitem[\citeproctext]{ref-imai2019should}
Imai, Kosuke, and In Song Kim. 2019. {``When Should We Use Unit Fixed
Effects Regression Models for Causal Inference with Longitudinal
Data?''} \emph{American Journal of Political Science} 63 (2): 467--90.

\bibitem[\citeproctext]{ref-jackson2004experience}
Jackson, Jonathan. 2004. {``Experience and Expression: Social and
Cultural Significance in the Fear of Crime.''} \emph{British Journal of
Criminology} 44 (6): 946--66.

\bibitem[\citeproctext]{ref-jha2026instrumental}
Jha, Akshaya, and Daniel Nagin. 2026. {``Instrumental Variables
Regression for Criminological Research.''} \emph{Annual Review of
Criminology} 9: 197--216.

\bibitem[\citeproctext]{ref-kim2023}
Kim, Jihoon, Lindsay Leban, Jun Sung Hong, and Yeungjeom Lee. 2023.
{``Sex Differences in the Overlap Between Bullying Perpetration and
Victimization: A Developmental Perspective.''} \emph{Journal of
Developmental and Life-Course Criminology} 9 (1): 35--53.
\url{https://doi.org/10.1007/s40865-022-00216-9}.

\bibitem[\citeproctext]{ref-king1994designing}
King, Gary, Robert O Keohane, and Sidney Verba. 1994. \emph{Designing
Social Inquiry: Scientific Inference in Qualitative Research}. Princeton
university press.

\bibitem[\citeproctext]{ref-lageson2012work}
Lageson, Sarah, and Christopher Uggen. 2012. {``How Work Affects
Crime---and Crime Affects Work---over the Life Course.''} In
\emph{Handbook of Life-Course Criminology: Emerging Trends and
Directions for Future Research}, 201--12. Springer.

\bibitem[\citeproctext]{ref-lanfear2024promise}
Lanfear, Charles C, and David S Kirk. 2024. {``The Promise and Perils of
the Sharing Economy: The Impact of Airbnb Lettings on Crime.''}
\emph{Criminology} 62 (4): 769--98.

\bibitem[\citeproctext]{ref-lanfear2020broken}
Lanfear, Charles C, Ross L Matsueda, and Lindsey R Beach. 2020.
{``Broken Windows, Informal Social Control, and Crime: Assessing
Causality in Empirical Studies.''} \emph{Annual Review of Criminology} 3
(1): 97--120.

\bibitem[\citeproctext]{ref-laub2006shared}
Laub, John H, and Robert J Sampson. 2006. \emph{Shared Beginnings,
Divergent Lives: Delinquent Boys to Age 70}. Harvard University Press.

\bibitem[\citeproctext]{ref-leszczensky2019deal}
Leszczensky, Lars, and Tobias Wolbring. 2019. {``How to Deal with
Reverse Causality Using Panel Data? Recommendations for Researchers
Based on a Simulation Study.''} \emph{Sociological Methods \& Research},
0049124119882473.

\bibitem[\citeproctext]{ref-llaudet2023data}
Llaudet, Elena, and Kosuke Imai. 2023. \emph{Data Analysis for Social
Science: A Friendly and Practical Introduction}. Princeton University
Press.

\bibitem[\citeproctext]{ref-lundberg2021your}
Lundberg, Ian, Rebecca Johnson, and Brandon M Stewart. 2021. {``What Is
Your Estimand? Defining the Target Quantity Connects Statistical
Evidence to Theory.''} \emph{American Sociological Review} 86 (3):
532--65.

\bibitem[\citeproctext]{ref-matsueda1998dynamics}
Matsueda, Ross L, and Kathleen Anderson. 1998. {``The Dynamics of
Delinquent Peers and Delinquent Behavior.''} \emph{Criminology} 36 (2):
269--308.

\bibitem[\citeproctext]{ref-mayo2018statistical}
Mayo, Deborah G. 2018. \emph{Statistical Inference as Severe Testing:
How to Get Beyond the Statistics Wars}. Cambridge University Press.

\bibitem[\citeproctext]{ref-moral2019dynamic}
Moral-Benito, Enrique, Paul D Allison, and Richard Williams. 2019.
{``Dynamic Panel Data Modelling Using Maximum Likelihood: An Alternative
to Arellano-Bond.''} \emph{Applied Economics} 51 (20): 2221--32.

\bibitem[\citeproctext]{ref-morgan2014counterfactuals}
Morgan, Stephen L, and Christopher Winship. 2014. \emph{Counterfactuals
and Causal Inference: Methods and Principles for Social Research}.
Cambridge University Press.

\bibitem[\citeproctext]{ref-mulder2021three}
Mulder, Jeroen D, and Ellen L Hamaker. 2021. {``Three Extensions of the
Random Intercept Cross-Lagged Panel Model.''} \emph{Structural Equation
Modeling: A Multidisciplinary Journal} 28 (4): 638--48.

\bibitem[\citeproctext]{ref-mulder2025time}
Mulder, Jeroen D, Manuel Voelkle, and Ellen L Hamaker. 2025. {``Time
Aggregation and Missing Time Frames in Causal Research with Panel
Data.''} \emph{PsyArXiv Preprings}.

\bibitem[\citeproctext]{ref-mundlak1978pooling}
Mundlak, Yair. 1978. {``On the Pooling of Time Series and Cross Section
Data.''} \emph{Econometrica: Journal of the Econometric Society},
69--85.

\bibitem[\citeproctext]{ref-muthuxe9n2024}
Muthén, Bengt, and Tihomir Asparouhov. 2024. {``Can Cross-Lagged Panel
Modeling Be Relied on to Establish Cross-Lagged Effects? The Case of
Contemporaneous and Reciprocal Effects.''} \emph{Psychological Methods},
May. \url{https://doi.org/10.1037/met0000661}.

\bibitem[\citeproctext]{ref-nagin2019real}
Nagin, Daniel S, and Robert J Sampson. 2019. {``The Real Gold Standard:
Measuring Counterfactual Worlds That Matter Most to Social Science and
Policy.''} \emph{Annual Review of Criminology} 2 (1): 123--45.

\bibitem[\citeproctext]{ref-nguyen2018measurement}
Nguyen, Holly, and Thomas A Loughran. 2018. {``On the Measurement and
Identification of Turning Points in Criminology.''} \emph{Annual Review
of Criminology} 1 (1): 335--58.

\bibitem[\citeproctext]{ref-nickell1981biases}
Nickell, Stephen. 1981. {``Biases in Dynamic Models with Fixed
Effects.''} \emph{Econometrica} 49 (6): 1417--26.

\bibitem[\citeproctext]{ref-nivette2026}
Nivette, Amy, and Manuel Eisner. 2026. {``Time Matters: Rethinking
Temporal Dynamics in Criminology.''} \emph{SocArXiv}.
\url{https://doi.org/10.31235/osf.io/ynx32_v1}.

\bibitem[\citeproctext]{ref-nivette2025assessing}
Nivette, Amy, and Isabelle van der Vegt. 2025. {``Assessing the Causal
Mechanisms Underlying Procedural Justice Theory in the Netherlands.''}

\bibitem[\citeproctext]{ref-obrien2019looking}
O'Brien, Daniel T, Chelsea Farrell, and Brandon C Welsh. 2019.
{``Looking Through Broken Windows: The Impact of Neighborhood Disorder
on Aggression and Fear of Crime Is an Artifact of Research Design.''}
\emph{Annual Review of Criminology} 2 (1): 53--71.

\bibitem[\citeproctext]{ref-oliveira2024aggressive}
Oliveira, Thiago R. 2024. {``Aggressive Policing and Undermined
Legitimacy: Assessing the Impact of Police Stops at Gunpoint on
Perceptions of Police in s{ã}o Paulo, Brazil.''} \emph{Journal of
Experimental Criminology} 20 (1): 83--121.

\bibitem[\citeproctext]{ref-oliveira2025}
---------. 2025. {``Legal Cynicism, Intrusive Policing, and the Dynamics
of Police Legitimacy: Evidence from Brazil{'}s Largest City.''}
\emph{Law \& Society Review} 59 (3): 548--83.
\url{https://doi.org/10.1017/lsr.2025.10}.

\bibitem[\citeproctext]{ref-ousey2011something}
Ousey, Graham C, Pamela Wilcox, and Bonnie S Fisher. 2011. {``Something
Old, Something New: Revisiting Competing Hypotheses of the
Victimization-Offending Relationship Among Adolescents.''} \emph{Journal
of Quantitative Criminology} 27 (1): 53--84.

\bibitem[\citeproctext]{ref-pager2003}
Pager, Devah. 2003. {``The Mark of a Criminal Record.''} \emph{American
Journal of Sociology} 108 (5): 937--75.
\url{https://doi.org/10.1086/374403}.

\bibitem[\citeproctext]{ref-pearl2009causality}
Pearl, Judea. 2009. \emph{Causality}. Cambridge University Press.

\bibitem[\citeproctext]{ref-pina2020reassessing}
Pina-Sánchez, Jose, and Ian Brunton-Smith. 2020. {``Reassessing the
Relationship Between Procedural Justice and Police Legitimacy.''}
\emph{Law and Human Behavior} 44 (5): 377.

\bibitem[\citeproctext]{ref-ramakers2017not}
Ramakers, Anke, Paul Nieuwbeerta, Johan Van Wilsem, and Anja Dirkzwager.
2017. {``Not Just Any Job Will Do: A Study on Employment Characteristics
and Recidivism Risks After Release.''} \emph{International Journal of
Offender Therapy and Comparative Criminology} 61 (16): 1795--1818.

\bibitem[\citeproctext]{ref-redcross2011more}
Redcross, Cindy, Megan Millenky, Timothy Rudd, and Valerie Levshin.
2011. {``More Than a Job: Final Results from the Evaluation of the
Center for Employment Opportunities (CEO) Transitional Jobs Program.''}
\emph{OPRE Report} 18.

\bibitem[\citeproctext]{ref-roman2026}
Roman, Gabriela Diana. 2026. {``The Method of Cascades: White Paper.''}
\emph{PsyArXiv}. \url{https://doi.org/10.31234/osf.io/y3unh}.

\bibitem[\citeproctext]{ref-rosenfeld2019}
Rosenfeld, Richard, and Joel Wallman. 2019. {``Did de{-}Policing Cause
the Increase in Homicide Rates?''} \emph{Criminology \& Public Policy}
18 (1): 51--75. \url{https://doi.org/10.1111/1745-9133.12414}.

\bibitem[\citeproctext]{ref-rubenstein2017causal}
Rubenstein, Paul K, Sebastian Weichwald, Stephan Bongers, Joris M Mooij,
Dominik Janzing, Moritz Grosse-Wentrup, and Bernhard Schölkopf. 2017.
{``Causal Consistency of Structural Equation Models.''} arXiv.
\url{https://doi.org/10.48550/arXiv.1707.00819}.

\bibitem[\citeproctext]{ref-rubin1986comment}
Rubin, Donald B. 1986. {``Comment: Which Ifs Have Causal Answers.''}
\emph{Journal of the American Statistical Association} 81 (396):
961--62.

\bibitem[\citeproctext]{ref-sampson1995crime}
Sampson, Robert J, and John H Laub. 1995. \emph{Crime in the Making:
Pathways and Turning Points Through Life}. Harvard University Press.

\bibitem[\citeproctext]{ref-sampson2016turning}
---------. 2016. {``Turning Points and the Future of Life-Course
Criminology: Reflections on the 1986 Criminal Careers Report.''}
\emph{Journal of Research in Crime and Delinquency} 53 (3): 321--35.

\bibitem[\citeproctext]{ref-sampson2006does}
Sampson, Robert J, John H Laub, and Christopher Wimer. 2006. {``Does
Marriage Reduce Crime? A Counterfactual Approach to Within-Individual
Causal Effects.''} \emph{Criminology} 44 (3): 465--508.

\bibitem[\citeproctext]{ref-sampson2004seeing}
Sampson, Robert J, and Stephen W Raudenbush. 2004. {``Seeing Disorder:
Neighborhood Stigma and the Social Construction of {`Broken
Windows'}.''} \emph{Social Psychology Quarterly} 67 (4): 319--42.

\bibitem[\citeproctext]{ref-skardhamar2014crime}
Skardhamar, Torbjørn, Christian W Monsbakken, and Torkild H Lyngstad.
2014. {``Crime and the Transition to Marriage: The Role of the Spouse's
Criminal Involvement.''} \emph{British Journal of Criminology} 54 (3):
411--27.

\bibitem[\citeproctext]{ref-skardhamar2014changes}
Skardhamar, Torbjørn, and Jukka Savolainen. 2014. {``Changes in Criminal
Offending Around the Time of Job Entry: A Study of Employment and
Desistance.''} \emph{Criminology} 52 (2): 263--91.

\bibitem[\citeproctext]{ref-sobel2008}
Sobel, Michael E. 2008. {``Identification of Causal Parameters in
Randomized Studies with Mediating Variables.''} \emph{Journal of
Educational and Behavioral Statistics} 33 (2): 230--51.
\url{https://doi.org/10.3102/1076998607307239}.

\bibitem[\citeproctext]{ref-solymosi2015}
Solymosi, Reka, Kate Bowers, and Taku Fujiyama. 2015. {``Mapping Fear of
Crime as a Context-Dependent Everyday Experience That Varies in Space
and Time.''} \emph{Legal and Criminological Psychology} 20 (2):
193--211. \url{https://doi.org/10.1111/lcrp.12076}.

\bibitem[\citeproctext]{ref-speyer2025importance}
Speyer, Lydia G, Xinxin Zhu, Yi Yang, Denis Ribeaud, and Manuel Eisner.
2025. {``On the Importance of Considering Concurrent Effects in
Random-Intercept Cross-Lagged Panel Modelling: Example Analysis of
Bullying and Internalising Problems.''} \emph{Multivariate Behavioral
Research} 60 (2): 328--44.

\bibitem[\citeproctext]{ref-stansfield2020}
Stansfield, Richard, Thomas J. Mowen, Laura Napolitano, and John H.
Boman. 2020. {``Examining Change in Family Conflict and Family Violence
After Release from Prison.''} \emph{Criminal Justice and Behavior} 47
(6): 668--87. \url{https://doi.org/10.1177/0093854820913311}.

\bibitem[\citeproctext]{ref-strotz1960recursive}
Strotz, Robert H, and H O A Wold. 1960. {``Recursive Vs. Nonrecursive
Systems: An Attempt at Synthesis (Part i of a Triptych on Causal Chain
Systems).''} \emph{Econometrica} 28 (2).

\bibitem[\citeproctext]{ref-sturgis2004using}
Sturgis, Patrick, Peter Smith, Ann Berrington, and Yiongjean Hu. 2004.
{``Using Repeated Measures Data to Analyse Reciprocal Effects: The Case
of Economic Perceptions and Economic Values.''} \emph{ESRC Methods
Festival}.

\bibitem[\citeproctext]{ref-taylor2001}
Taylor, Ralph B. 2001. \emph{Breaking Away from Broken Windows:
Baltimore Neighborhoods and the Nationwide Fight Against Crime, Grime,
Fear, and Decline}. Crime \& Society. Boulder, Colo: Westview Press.

\bibitem[\citeproctext]{ref-uggen2008have}
Uggen, Chris, and Sara Wakefield. 2008. {``What Have We Learned from
Longitudinal Studies of Work and Crime?, w: The Long View of Crime: A
Synthesis of Longitudinal Research, Red.''} \emph{AM Liberman, New
York}.

\bibitem[\citeproctext]{ref-unnever2020}
Unnever, James D., and Cecilia Chouhy. 2020. {``Race, Gender, and
Perceptions of Peer Delinquency: A Within-Subject Analysis.''}
\emph{Deviant Behavior} 41 (11): 1413--30.
\url{https://doi.org/10.1080/01639625.2019.1623604}.

\bibitem[\citeproctext]{ref-vaisey2017you}
Vaisey, Stephen, and Andrew Miles. 2017. {``What You Can---and
Can't---Do with Three-Wave Panel Data.''} \emph{Sociological Methods \&
Research} 46 (1): 44--67.

\bibitem[\citeproctext]{ref-voelkle2018role}
Voelkle, Manuel C, Christian Gische, Charles C Driver, and Ulman
Lindenberger. 2018. {``The Role of Time in the Quest for Understanding
Psychological Mechanisms.''} \emph{Multivariate Behavioral Research} 53
(6): 782--805.

\bibitem[\citeproctext]{ref-wiesner2023}
Wiesner, Margit, Deborah M. Capaldi, David C. R. Kerr, and Weiwei Wu.
2023. {``Bidirectional Associations of Mental Health with Self-Reported
Criminal Offending over Time for at-Risk Early Adult Men in the USA.''}
\emph{Journal of Developmental and Life-Course Criminology} 9 (2):
211--37. \url{https://doi.org/10.1007/s40865-022-00221-y}.

\bibitem[\citeproctext]{ref-wikstruxf6m2024}
Wikström, Per-Olof H., Kyle Treiber, and Gabriela D. Roman. 2024.
\emph{Character, circumstances, and criminal careers: towards a dynamic
developmental and life-course criminology}. First edition. Clarendon
studies in criminology. Oxford: Oxford University Press.
\url{https://doi.org/10.1093/oso/9780198865865.001.0001}.

\bibitem[\citeproctext]{ref-kelling1982broken}
Wilson, James Q, and George L Kelling. 1982. {``Broken Windows.''}
\emph{Atlantic Monthly} 249 (3): 29--38.

\bibitem[\citeproctext]{ref-wooldridge2010econometric}
Wooldridge, Jeffrey M. 2010. \emph{Econometric Analysis of Cross Section
and Panel Data}. MIT press.

\bibitem[\citeproctext]{ref-zyphur2020data}
Zyphur, Michael J, Paul D Allison, Louis Tay, Manuel C Voelkle,
Kristopher J Preacher, Zhen Zhang, Ellen L Hamaker, et al. 2020. {``From
Data to Causes i: Building a General Cross-Lagged Panel Model (GCLM).''}
\emph{Organizational Research Methods} 23 (4): 651--87.

\end{CSLReferences}




\end{document}
